\section{Bounds on intervals and the critical line}
\label{sec:zeros}

\subsection{One family for the lattice and Fourier estimates}
\label{sec:assembly}

Take the conductor-one target \(\psi=1\). Fix the loss exponents
\(e_D=2,\ e_L=0\) and the balance constant \(C_0=3\) used in the
\releasecite{Development/AutomorphicArithmetic/ThetaCorrespondence/ThetaPairLifts/ThirdPlaceSeries/TopShape.lean}{50}{window construction}.

\teamanchor{proof.join-size-requirements}%
The uniform comparisons of Sections~\ref{sec:family}--\ref{sec:fourier}
have finitely many precision cutoffs and size requirements. A lower
bound may be a function of earlier constants. At fixed precision take
the pointwise maximum \(K_*\) of the lower bounds in each position.
In the position for \(C_H\) include the bound
\(1+\max\{4P_t+1,9P_t+2,9P_t+1,C_{F,0}\}\)
from \eqref{eq:raw-height-cap}. Here \(P_t\) comes from the twist
cutoff and \(C_{F,0}=c_F+1\) comes from the row height after
the finite level change. Both are functions of the earlier data.
Thus the same \(C_H\) is available for the height and the trace
accuracy before choosing a scale or translate.
A family meeting \(K_*\) therefore meets all of the size
requirements. Each comparison also requires its hypotheses on the
local models and maps.
\releasecite{Formalism/SharpRestriction/ParamsHeightRankConstant.lean}{138}{Staged lower bounds}

\medskip
\noindent
\teamanchor{proof.choose-compatible-family}%
\textbf{Simultaneous family selection (input).}
Assume the availability of the arithmetic hosts on intervals from
Section~\ref{sec:family}. There is an existence cutoff
\(\varepsilon_{\rm ex}>0\). For each
\(0<\varepsilon<\min(1/52,\varepsilon_{\rm ex})\) and each
system \(K\) of staged lower bounds there exist one family
\(\mathfrak F\) and one exponent \(C_H\ge0\) meeting \(K\).
The family has the following data and conditions.

\begin{enumerate}
\item The set \(\operatorname{Win}(\mathfrak F)\) of real scales
contains intervals \([Y,Y^{3/2}]\) with arbitrarily large \(Y\).
Each scale in this set has a nonempty index set \(I_x\). Beyond a
family threshold the indices use the prescribed arithmetic hosts
at \(x\).

\item At each such scale an index specifies \(E/F\), a coefficient
field with its places, the ranks and weights of the towers, the
levels, the words and their local twists, the period spaces with
their row and column coordinates, and the corresponding test and
integral lattice data. A word, tuple, stage and place specify its
local Hecke data and intended test in these coordinates.
At a fixed \(x\), the indices have specified identifications
of the ranks, word labels and their letters, \(q,d,Q_S\), the
Artin coordinates, the root discriminant, the split primes
\(g_0,g_1\), the two further level primes \(p,p'\), the uniform family
constants
and \(T,T_2\). The rank, level-prime and
sample cardinality conditions are those in Section~\ref{sec:family}.
The stage weights and the range for \(b_t\) are the raw intervals
in Section~\ref{sec:fourier-sampling}.

With \(b=\exp(x^\varepsilon)\), the level of each word satisfies
\(b^{m/2}\le T_w\) and \(\log T_w\le C_T m\log b\).
Write \(f=[F:\mathbf Q]\) and \(R_m=k+m(t+1)=\operatorname{rank}X_m\).
The positive integer \(D_\Gamma\) at this index is a common
multiple of the class stabilizer orders in the column chains.
Its prime divisors are at most \(2fR_m+1\) and
\(\log D_\Gamma\le C_\Delta fR_m\log(fR_m)\).
Here \(C_T,C_\Delta\) are uniform family constants. These bounds on
\(T_w\) and \(D_\Gamma\) are derived below, in the paragraphs on the
size of the levels and on the stabilizer multiple. The first uses the
column data of the record, as stated there.
The column dimensions satisfy \eqref{eq:trace-balance} and
\[
 C_0\Gamma^2x^{C_H}\le |W|\le C_0\Gamma^2x^{C_H}+1
\]
in the trace notation, with \(C_0=3\). Only the lower bound is used
below. The balance \eqref{eq:trace-balance} is not a separate choice.
Section~\ref{sec:trace} derives it from the word-prime interval of the
level choice and the level index.

The selection also fixes two exponents \(u_{\rm cl},u_{\rm node}\),
taken to be at least the cutoffs of the eventual estimates, before any
scale. Put \(\Delta_{\rm cl}=\lceil x^{u_{\rm cl}}\rceil\)
and \(H_{\rm node}=2\lceil x^{u_{\rm node}}\rceil\). Every index at
\(x\) carries, for every word, the Harish--Chandra rows of
Section~\ref{sec:integral} with these spacings. Consecutive entries in
the initial block differ by \(4\Delta_{\rm cl}\), and the later nodes
are \(B_\tau+lH_{\rm node}\) and \(B_\tau+(l+\tfrac12)H_{\rm node}\)
\releasecite{Development/AbelianVarietyTheory/Definitions/GlobalTwisting/PairRowIffQdRows.lean}{93}{Prescribed row values}.
Every index also carries its own embedded geometric coefficient field
\(K_0\), fixed with the base column choice before the stages
\releasecite{Development/AbelianVarietyTheory/Definitions/AutomorphicMembership/PositionOfIntro.lean}{23}{Field fixed with the base column choice}.
Both auxiliary levels are admissible for \(K_0\)
\releasecite{Development/AbelianVarietyTheory/Proofs/IwasawaDecomposition/IndexBoundCoefficientFieldOfExactSmallField.lean}{28}{Admissibility at both levels},
and \(s(K_0)\le P_O\) beyond the family threshold
\releasecite{Development/AbelianVarietyTheory/Proofs/HeightProofs/EventuallyConstructedPointHostSizeNumerical.lean}{19}{Size of the field}.

\item The record includes admissible integers \(a_t,b_t\) on \(T\).
In particular \(a_E\le a_t\le a_E+a_E/(kx^{13})\), where \(a_E=xb\).
For any seed and any \(u\in H\) there is an index whose translate
is \(u\) and whose two arrays agree with the seed's under the common
labeling of \(T\). The weights at \(u\) are read from the stage data
of this translated index.
\end{enumerate}

This statement uses the selection part of the release's
\releasecite{Development/AutomorphicArithmetic/ThetaCorrespondence/ThetaPairLifts/ThirdPlaceSeries/LinkEx.lean}{15}{design existence theorem}
and the \releasecite{Formalism/SharpRestriction/FamCore.lean}{76}{common data and the translate property}.
In the release the existence of the design choices is a hypothesis of
that theorem. It is discharged by the
\releasecite{Development/AutomorphicArithmetic/CharacterBoxes/ClassSel.lean}{14}{class selection},
which builds the hosts, ranks, level primes, word sets and tied kits
scale by scale and then chooses constants above the staged lower
bounds. This edition reconstructs the host sizes and rank arithmetic
in Section~\ref{sec:family}, the column balance in
Section~\ref{sec:trace} and the joining of the size requirements at
the start of this subsection. The paragraphs below give the order of
the construction and derive its analytic and combinatorial inputs:
the level primes \(T_w\) and their size, the stabilizer multiple
\(D_\Gamma\), the windows for the arrays, the split primes \(g_0,g_1\)
with disjoint bad sets, the primes \(p,p'\), the word sets, and the ties and the
translate property of item~3. The existence of the constructed record
at each translate is taken from the release. The residual is stated
precisely at the end of these paragraphs.
The character estimates and the Selmer input enter through the
arithmetic hosts of Section~\ref{sec:family}. The local, trace and
height comparisons enter through separate hypotheses in the
\releasecite{Development/AutomorphicArithmetic/CharacterBoxes/DesWin.lean}{11}{subsequent assembly in the release}.

\paragraph{Order of the construction.}
\teamanchor{proof.selection-construction-order}%
The release makes the choices in the following order. Each threshold
depends only on the data fixed before it.
\begin{enumerate}
\item[(a)] \emph{Precision and constants.} The existence cutoff
satisfies \(\varepsilon_{\rm ex}\le1/53\)
\releasecite{Development/AutomorphicArithmetic/CharacterBoxes/DesignS.lean}{21}{Cutoff and gluing at one precision}.
Fix \(\varepsilon\). The host constants, including \(C_8\), come
with the host statement at this precision. The remaining uniform
family constants, including \(C_\Gamma\), \(C_H\) and \(C_T\), are
then chosen in the staged order at or above the pointwise maxima
\(K_*\) described above
\releasecite{Development/AbelianVarietyTheory/Proofs/GlobalInterpolation/ExistsStdLarge.lean}{39}{Constants above the staged lower bounds}.
After them come the exponents \(u_{\rm cl},u_{\rm node}\). This step
only takes finitely many maxima.
\item[(b)] \emph{Scales, ranks and host.} The statement assumes the
hosts on intervals of Section~\ref{sec:family}. They supply, at each
scale \(x\) beyond a threshold in a union of intervals \([Y,Y^{3/2}]\)
with \(Y\) unbounded, the ranks and the host data
\(q,d,Q_S,g_0,g_1,p,p',q_0,T,T_2\). The existence of \(g_0,g_1\) and
\(p,p'\) is derived below.
\item[(c)] \emph{Arrays.} Integers \(a_t,b_t\) on \(T\) are chosen in
the raw windows. In the release they come with the host and satisfy a
further admissibility condition
\releasecite{Development/AbelianVarietyTheory/Definitions/GlobalTwisting/AuxiliaryPrimeMemAvoidedPrimes.lean}{122}{Hosts with admissible arrays}.
\teamanchor{proof.selection-array-windows}%
Both windows are nonempty. For \(a_t\) the window
\([a_E,a_E+a_E/(kx^{13})]\) has length at least \(1\), because
\(a_E=xb\), \(k\le x^5\) and \(b=\exp(x^\varepsilon)\ge x^{17}\) at
large \(x\). The window for \(b_t\) is
\([b',b'(1+2n\delta_1x^\varepsilon L_{\mathrm{a},G}/k)]\) in the
notation of Section~\ref{sec:fourier-sampling}. Its length is
\(2n\,\delta_1x^\varepsilon L_{\mathrm{a},G}\,b'/k\). Here \(n\ge1\),
\(\delta_1x^\varepsilon=x^{7\varepsilon/8}\ge1\), and
\(b'=x\exp(x^\varepsilon-x^{\varepsilon/3})\ge x^5\ge k\) at large
\(x\). The parameter is \(L_{\mathrm{a},G}=C_{\mathrm a}N_1^2\log a_E\)
with a family constant \(C_{\mathrm a}\ge1\)
\releasecite{Formalism/SharpRestriction/GrowthBounds/IsResidueClassZero.lean}{58}{The window parameter},
and the release shows \(N_1\ge1\) and
\(x^\varepsilon\le L_{\mathrm{a},G}\) beyond a threshold depending
only on \(\varepsilon\)
\releasecite{Development/AbelianVarietyTheory/Proofs/GlobalInterpolation/GoodTopCountLeConstMulLog.lean}{65}{Lower bound for the window parameter}.
The same statement gives
\(L_{\mathrm{a},G}\le2C_{\mathrm a}C_N^2(\log x)^2x^\varepsilon\) with a
constant \(C_N\) depending only on \(\varepsilon\), and this is at most
\(x^{\varepsilon+\varepsilon/100}\) at large \(x\).
So the length is at least \(1\), and for instance the constant arrays
\(a_t=\lceil a_E\rceil\), \(b_t=\lceil b'\rceil\) lie in the windows.
This is the release's
\releasecite{Development/AbelianVarietyTheory/Proofs/GlobalInterpolation/ExistsWeightsInWindows.lean}{18}{Arrays in the raw windows}.
The arrays carried by the records are chosen again inside these
windows so that they also satisfy the admissibility condition. This
choice is (R3) below.
\item[(d)] \emph{Words.} A word set \(W\) of the required size is
chosen as below.
\item[(e)] \emph{Records at the translates.} For each \(u\in H\) one
record is constructed from the same ranks, host, arrays \(a,b\) and
words, with translate \(u\). Inside this construction the release
first fixes the column data \(d_w^{(0)},D_\star\). Then the level
primes \(T_w\) are chosen in the windows of the next paragraph. The
record over these levels is built last
\releasecite{Development/AbelianVarietyTheory/Proofs/GlobalInterpolation/ExistsObjectsOfAdmissibleTwistDatumStatement.lean}{8}{Level primes inside the record construction}.
\item[(f)] \emph{Ties and translates.} The index set at \(x\) is the
set of translates. All records at \(x\) are built from identical
ranks, host data, arrays and words, so the identifications at fixed
\(x\) in item~2 hold. For a seed \(i_0\) and \(u\in H\), the index
\(i(u)=u\) has translate \(u\) and the arrays of \(i_0\) on the common
labeling of \(T\). This is item~3. The release records it in the
\releasecite{Development/AutomorphicArithmetic/CharacterBoxes/KitS.lean}{14}{Kit at one scale}.
\item[(g)] \emph{Gluing.} The kits over the window scales are glued
into one family with the common constants. The family threshold is
the maximum of the thresholds in (b)--(f)
\releasecite{Development/AbelianVarietyTheory/Proofs/GlobalInterpolation/ExistsFamPOfKits.lean}{159}{Gluing the kits}.
\end{enumerate}

\paragraph{Split primes in short intervals.}
\teamanchor{proof.level-primes-short-intervals}%
Let \(d>3\) be a prime with \(d\equiv3\bmod4\) and put
\(K_d=\mathbf Q(\sqrt{-d})\). Let \(M\ge3\) be an integer and put
\(B=4dM^2\). Call a rational prime \(T\) an \emph{\(M\)-norm prime}
if \(T=\pi\overline\pi\) with \(\pi\in\mathcal O_{K_d}\) and
\(\pi\equiv1\bmod M\mathcal O_{K_d}\). Then also
\(\overline\pi\equiv1\), so \(T-1\in M\mathcal O_{K_d}\cap\mathbf Z
=M\mathbf Z\) and \(T\equiv1\bmod M\). We show the following. There are
an absolute effective constant \(c\in(0,1]\) and an absolute \(X_0\)
such that, if \(X\ge X_0\), \(\log B\le c\sqrt{\log X}\) and
\(e^{-c\sqrt{\log X}}\le h\le1\), then
\begin{equation}
\label{eq:selection-short-norm-primes}
 \#\{T\in(X,X(1+h)]:\ T\text{ is an }M\text{-norm prime}\}
 \ \ge\ \frac{hX}{4B\log(2X)} .
\end{equation}

Let \(L\) be the ray class field of \(K_d\) modulo
\(\mathfrak m=M\mathcal O_{K_d}\). Since \(\mathfrak m\) is stable
under complex conjugation, \(L/\mathbf Q\) is Galois. Put
\(G=\operatorname{Gal}(L/\mathbf Q)\), \(H_L=\operatorname{Gal}(L/K_d)\)
and \(C=\{1\}\). By the Artin isomorphism of the ray class field
\cite[Ch.~V, Theorem~3.6 and the paragraph after it]{f1milne-cft}, a prime ideal
\(\mathfrak p\nmid\mathfrak m\) of \(K_d\) splits completely in \(L\)
exactly when it lies in the trivial ray class, that is, when
\(\mathfrak p=(\pi)\) with \(\pi\equiv1\bmod\mathfrak m\). There is no
sign condition because \(K_d\) has no real place. Hence every
rational prime \(T\) that splits completely in \(L\) is an \(M\)-norm
prime. The group \(H_L\) is abelian and contains the class \(C\). Its
characters are ray class characters modulo \(\mathfrak m\), so the
quantity \(\mathcal Q\) of \cite[arXiv v3, (1.11)]{s1-thorner-zaman-chebotarev}
satisfies \(\mathcal Q\le N\mathfrak m=M^2\). Hence
\(D_{K_d}\mathcal Q\,n_{K_d}^{n_{K_d}}\le4dM^2=B\). The number of reduced
forms of discriminant \(-d\) is at most \(d\), so \(h(K_d)\le d\).
Also \(|\mathrm{Cl}_{\mathfrak m}(K_d)|\le h(K_d)N\mathfrak m\). Thus
\(|G|=2|\mathrm{Cl}_{\mathfrak m}(K_d)|\le2dM^2\le B\).

Apply \cite[arXiv v3, Theorem~1.4]{s1-thorner-zaman-chebotarev} with
\(F=\mathbf Q\), the abelian subgroup \(H_L\) and \(g_C=1\). The error
there increases with \(\log(D_{K_d}\mathcal Q n_{K_d}^{n_{K_d}})\), so we
may replace that quantity by \(\log B\). For \(x\ge B^{c_2}\) this gives
\[
 \pi_C(x)=\frac1{|G|}\bigl(\operatorname{Li}(x)-\theta_1\operatorname{Li}(x^{\beta_1})\bigr)
 \bigl(1+O(\mathcal E(x))\bigr),\qquad
 \mathcal E(x)=e^{-c_3\log x/\log B}+e^{-(c_3\log x/2)^{1/2}},
\]
where \(\theta_1=\chi_1(1)=1\) if the exceptional zero \(\beta_1\)
exists and \(\theta_1=0\) otherwise. Here \(c_2,c_3\) and the implied
constant are absolute and effective
\cite[arXiv v3, Section~2]{s1-thorner-zaman-chebotarev}.
Suppose \(\beta_1\) exists. It is a zero of the Hecke \(L\)-function
of a real character \(\chi_1\) of \(H_L\), and
\(1-\beta_1<1/(8\log(4d\mathcal Q))\). Let \(K'\) be the field cut out
by \(\chi_1\). Thus \(K'=K_d\) if \(\chi_1\) is trivial, and otherwise
\(K'/K_d\) is quadratic. Then \(\zeta_{K'}\) is \(\zeta_{K_d}\) or
\(\zeta_{K_d}L(s,\chi_1)\), so \(\zeta_{K'}(\beta_1)=0\). Moreover
\([K':\mathbf Q]\le4\), and by the conductor--discriminant formula
\(D_{K'}\le d^2N\mathfrak f_{\chi_1}\le\min\{B^2,(4d\mathcal Q)^2\}\).
Thus \(1-\beta_1<1/(4\log D_{K'})\), and \(\beta_1\) is the possible
Landau--Siegel zero of \(\zeta_{K'}\). Stark's bound, as recorded in
\cite[arXiv v3, (1.4)]{s1-thorner-zaman-chebotarev} for the field
\(K'\), gives \(1-\beta_1\gg(n_{K'}^{n_{K'}}\log D_{K'}+D_{K'}^{1/n_{K'}})^{-1}\).
Here \(n_{K'}\in\{2,4\}\), and \(D_{K'}^{1/n_{K'}}\le B^{1/2}\) in both
cases, since \(D_{K_d}=d\le B\) and \(D_{K'}\le B^2\). Hence
\(1-\beta_1\ge c_SB^{-1/2}\) with an absolute effective \(c_S>0\).

Take \(c=\min\{1,(c_3/8)^{1/2}\}\). Then
\(c_3\log X/\log B\ge(c_3/c)\sqrt{\log X}\ge2c\sqrt{\log X}\) and
\((c_3\log X/2)^{1/2}\ge2c\sqrt{\log X}\). Hence
\(\mathcal E(X)\le2e^{-2c\sqrt{\log X}}\le2he^{-c\sqrt{\log X}}\), and
\(\mathcal E\) decreases in \(x\). Since \(\log B\le c\sqrt{\log X}\),
we have \(X\ge B^{c_2}\) for \(X\ge X_0\). Also
\((1-\beta_1)\log X\ge c_S e^{-(c/2)\sqrt{\log X}}\log X\ge\log2\), so
\(X^{\beta_1-1}\le1/2\). For \(t\ge X\) the derivative of the main
term \(\operatorname{Li}(t)-\theta_1\operatorname{Li}(t^{\beta_1})\) is
\((1-\theta_1t^{\beta_1-1})/\log t\ge1/(2\log t)\). Hence the main term
increases on \((X,X(1+h)]\) by at least \(hX/(2|G|\log(2X))\).
At either endpoint the main term lies between \(0\) and
\(\operatorname{Li}(2X)\le4X/\log X\). The two error terms therefore
total at most \(16C_{\rm TZ}hXe^{-c\sqrt{\log X}}/(|G|\log X)\), where
\(C_{\rm TZ}\) is the implied constant. For \(X\ge X_0\) this is at
most \(hX/(4|G|\log(2X))\). Subtracting gives at least
\(hX/(4|G|\log(2X))\) primes in \((X,X(1+h)]\) that split completely
in \(L\). All of them are \(M\)-norm primes, and \(|G|\le B\). This
proves \eqref{eq:selection-short-norm-primes}.
The release proves a statement of the same form, with \(B\) replaced
by a power \(B^A\) and with \(d\ge100\)
\releasecite{Development/AbelianVarietyTheory/Proofs/GlobalInterpolation/NormPrimesImagQuadShortIntervalLowerBound.lean}{98}{Norm primes in short intervals}.

\emph{Level primes.} At a record over the host, take \(M=Q_Sqd(g_0g_1)^2\),
\(X=Y_w\) with \(Y_w\) as in Section~\ref{sec:trace}, and
\(h=1/(fk)\). By the size of the levels derived below,
\(\log Y_w\ge c_sm x^\varepsilon\) with \(c_s=14\). Section~\ref{sec:family}
gives \(m\ge x^{0.49}\), \(f\le k\le x^5\),
\(\log(Q_Sqd)\le2x^{\varepsilon/5}(\log x)^2\) and
\(\log g_i\le9x^\varepsilon\). Hence \(\log M\le37x^\varepsilon\) and
\(\log B\le120x^\varepsilon\) at large \(x\). Since \(\varepsilon<1/52\),
\(c\sqrt{\log Y_w}\ge c\,c_s^{1/2}x^{0.245}\) exceeds both
\(120x^\varepsilon\) and \(10\log x\ge\log(fk)\) at large \(x\).
Thus \eqref{eq:selection-short-norm-primes} gives at least
\(\exp(\log Y_w-10\log x-120x^\varepsilon-\log(4\log2Y_w))>1\) primes
\(T_w\) in \([Y_w,Y_w(1+1/(fk))]\). Each has a norm generator
\(\pi_{T_w}\equiv1\bmod Q_Sqd(g_0g_1)^2\), so
\(T_w\equiv1\bmod Q_Sqd(g_0g_1)^2\). As for the five host primes in
Section~\ref{sec:family}, it follows that \(T_w\) splits completely in
\(E\) and \(F\). Also \(T_w>Y_w>\max\{p,p'\}\), since
\(\log p,\log p'\le C_8\log(Q_Sqd)\) is smaller than \(c_smx^\varepsilon\).
This is the word-prime interval of Section~\ref{sec:trace}.

\emph{Size of the levels.}
\teamanchor{proof.selection-level-bounds}%
Inside the record construction (e), the column data are fixed before
the levels. By the index formula of Section~\ref{sec:trace},
\(d_w^{(0)}=d_w^*\), and \(d_w^*\) and \(D_\star\) are the
mass--Vandermonde products of Section~\ref{sec:integral}. These
products do not involve \(T_w\). In the notation of
\eqref{eq:arithmetic-level-quotient} put
\(X_w=(D_\star/d_w^*)^{1/(fk)}\), so that
\(Y_w=X_w\Gamma^{-2/(fk)}\), and let \(L=x^\varepsilon=\log b\).
We bound the terms of \eqref{eq:arithmetic-level-quotient} using the
estimates proved after it, whose implied constants are absolute.
\begin{itemize}
\item \(0\le\log(v_0/\pi)\le\log(2v_0)\le L\). Since \(m\le k\), the
coefficient of \(\log(v_0/\pi)\) lies in \([0,2m]\), and the
\(\log(2v_0)\) term has modulus at most \(L\).
\item The split primes lie in \([(kb)^8,2(kb)^8]\)
(Section~\ref{sec:family}), and
\(\log2+8\log k\le L\) at large \(x\). So
\(16L\le\lambda_g\le18L\), and the coefficient of \(\lambda_g\) lies
in \([m,2m]\).
\item At each of the \(f\) embeddings \(\rho\), the log Vandermonde
ratio is at least \(k\sum_j\log w_j-O(kL)\) and at most \(O(kmL)\).
Every stage weight is at least \(1\)
(Section~\ref{sec:fourier-sampling}). After division by \(fk\) the
Vandermonde term therefore lies between \(-O(L)\) and \(O(mL)\).
\end{itemize}
Hence \(\log X_w\ge16mL-O(L)\ge15mL\) once \(m\ge x^{0.49}\) is large,
and \(\log X_w\le C_XmL\) for an absolute constant \(C_X\). Next,
\(\log\Gamma=C_\Gamma m^5f\log b\) and \(m^5f\log b\le128kx^{-21\varepsilon/16}\)
give
\[
 0\le\frac{2\log\Gamma}{fk}\le\frac{256\,C_\Gamma}{f}\,x^{-21\varepsilon/16}\le1
\]
beyond a threshold that depends only on \(C_\Gamma\) and \(\varepsilon\).
Thus
\[
 14\,m\log b\le\log X_w-1\le\log Y_w\le\log X_w\le C_Xm\log b .
\]
For \(T_w\in[Y_w,Y_w(1+1/(fk))]\) this gives \(T_w\ge Y_w\ge b^{14m}\ge b^{m/2}\)
and \(\log T_w\le\log Y_w+1\le(C_X+1)m\log b\). So item~2 holds with
any \(C_T\ge C_X+1\). This is one more staged lower bound in (a), and
it involves only an absolute constant. The bounds agree with
\eqref{eq:arithmetic-count-at}. They are derived here without the
record (R1). The edition still takes from the release the
identification of the column data with the mass--Vandermonde products
and the grid conditions, as cited in Section~\ref{sec:integral}
\releasecite{Development/AbelianVarietyTheory/Proofs/FlagMeasures/ExactRanksOfRecord.lean}{11}{Column coefficient and row main term}
\releasecite{Formalism/SharpRestriction/GlobalSchurBounds/StandardLattice.lean}{500}{Grid conditions}.

\emph{The stabilizer multiple.}
\teamanchor{proof.selection-stabilizer-multiple}%
Take \(D_\Gamma=\operatorname{lcm}_c|\Gamma_c|\) over the class
stabilizers \(\Gamma_c\) in the column chains, as in
Section~\ref{sec:integral}. Each \(\Gamma_c=G(F)\cap cJc^{-1}\) is taken in the unitary group
\(G\) of one of the definite spaces \(X_j\) of the column tower, of
rank \(k+j(t+1)\le R_m\) over \(E\). The open compact level \(J\)
fixes an \(\mathcal O_E\)-lattice \(\Lambda\), for instance the sum of
the finitely many translates of a standard lattice. So \(\Gamma_c\)
preserves the lattice \(c\Lambda\) in \(X_j\). It is finite by
definiteness, and it acts faithfully on this lattice, a free
\(\mathbf Z\)-module of rank \(n\le2fR_m\). So \(\Gamma_c\) is a finite
subgroup of \(\mathrm{GL}_n(\mathbf Z)\subset\mathrm{GL}_n(\mathbf Q)\). By Minkowski's
theorem \cite[Theorem~1]{s2a-guralnick-lorenz}, the order of every
finite subgroup of \(\mathrm{GL}_n(\mathbf Q)\) divides
\[
 M(n)=\prod_\ell\ell^{e_\ell(n)},\qquad
 e_\ell(n)=\sum_{i\ge0}\left\lfloor\frac{n}{\ell^i(\ell-1)}\right\rfloor .
\]
Each \(e_\ell(n)\) is nondecreasing in \(n\), so \(D_\Gamma\) divides
\(M(2fR_m)\). Moreover \(e_\ell(n)=0\) for \(\ell>n+1\), so every
prime divisor of \(D_\Gamma\) is at most \(2fR_m+1\). For the size,
\(e_\ell(n)\le n\ell/(\ell-1)^2\le4n/\ell\). For an integer
\(N\ge1\),
\[
 N\log N\ge\log N!\ge\sum_{\ell\le N}\Bigl(\frac N\ell-1\Bigr)\log\ell
 \ge N\sum_{\ell\le N}\frac{\log\ell}{\ell}-N\log4 .
\]
Here we used \(\prod_{\ell\le N}\ell<4^N\). This gives
\(\sum_{\ell\le N}\log\ell/\ell\le\log N+\log4\). With \(N=n+1\),
\[
 \log D_\Gamma\le\log M(n)\le4n\bigl(\log(n+1)+\log4\bigr).
\]
For \(n=2fR_m\) and \(fR_m\ge2\) the right side is at most
\(8fR_m(\log(fR_m)+\log12)\le40fR_m\log(fR_m)\). So item~2 holds with
\(C_\Delta=40\). This is the Minkowski integer used for the
stabilizer allowance in Section~\ref{sec:integral}.

\emph{The split primes \(g_0,g_1\).} Take \(M=Q_Sqd\), \(X=(kb)^8\) and
\(h=1\). Then \(\log B\le7x^{\varepsilon/5}(\log x)^2\), which is at most
\(c(8x^\varepsilon)^{1/2}\le c\sqrt{\log X}\) at large \(x\). So
\eqref{eq:selection-short-norm-primes} gives at least
\(X/(4B\log2X)\ge X^{9/10}\) primes
\(g\in((kb)^8,2(kb)^8]\) with \(g=\pi\overline\pi\) and
\(\pi\equiv1\bmod Q\mathcal O_{K_d}\). These are the norm conditions of
Section~\ref{sec:family}. For the disjointness of the bad sets
\eqref{eq:classical-level-bad-primes}, fix one such \(g_0\) and put
\(S=B_r(g_0)\). Every \(\ell\in S\setminus\{g_0\}\) divides some
\(g_0^a-1\) with \(a\le r\). Hence
\(|S|\le1+\sum_{a\le r}a\log_2(2X)\le r^2\log(2X)\), since \(r\ge k\)
is large. A candidate \(g\) satisfies \(B_r(g)\cap S\ne\varnothing\)
exactly when \(g\in S\), or when there are \(\ell\in S\) and \(a\le r\)
with \(\ell^{4v_\ell(g^a-1)}>g^a\ge X^a\). In the second case
\(g^a\equiv1\) modulo the least power \(\ell^{v}>X^{a/4}\). The
congruence \(y^a\equiv1\bmod\ell^v\) has at most \(2a\) solutions,
and each class contains at most \(2X/\ell^v+1\le3X^{3/4}\) integers in
\((X,2X]\). So at most
\(|S|(1+\sum_{a\le r}6aX^{3/4})\le7r^4X^{3/4}\log(2X)\) candidates are
excluded. With \(r\le x^5\) and \(X\ge\exp(8x^\varepsilon)\), this is
less than \(X^{9/10}-1\) at large \(x\). Hence some \(g_1\ne g_0\) has
\(B_r(g_0)\cap B_r(g_1)=\varnothing\). The release proves the same
count
\releasecite{Development/AbelianVarietyTheory/Proofs/GlobalInterpolation/ExistsTwoLargeNormOnePrimes.lean}{54}{Two split primes with disjoint bad sets}.

\emph{The primes \(p,p'\).} Put \(m_1=\operatorname{lcm}(15,Q)\le15Q\)
and \(u=m_1^A\), with a fixed \(A\ge12\) as large as in the argument for
\(q_0\) in Section~\ref{sec:family}. By
\cite[arXiv v2, Corollary~1.4]{f1thorner-zaman}, the sum of \(\log p\)
over the primes \(p\le u\) with \(p\equiv1\bmod m_1\) equals
\(\lambda u/\varphi(m_1)\) times \(1+O(e^{-c_4A}+e^{-c_4(\log u)^{3/5}(\log\log u)^{-1/5}})\).
For fixed large \(A\) and large \(x\) this relative error has modulus
less than \(1/2\). Here \(\lambda=1\) if there is no exceptional zero
modulo \(m_1\). Otherwise
\(\lambda>\tfrac18\min\{1,(1-\beta_{m_1})\log u\}\)
\cite[arXiv v2, (1.7)]{f1thorner-zaman}, and
\(1-\beta_{m_1}\ge b_0m_1^{-1/2}\) as recalled in
Section~\ref{sec:family}. Since \(\log u\ge1\), in both cases
\(\lambda\ge\tfrac18b_*m_1^{-1/2}\) with \(b_*=\min\{1,b_0\}\).
Hence the sum is at least \(b_*m_1^{A-3/2}/16\), which exceeds
\(\log u=A\log m_1\) at large \(x\). One prime contributes at most
\(\log u\). So there are two distinct primes
\(p,p'\equiv1\bmod m_1\) with \(p,p'\le m_1^A\le Q^{2A}\), since
\(Q\ge15\). So \(C_8=2A\) suffices.

\paragraph{Word sets.}
\teamanchor{proof.word-sets}%
The letters of a word are \(k-2\) bits, and a word set is an injective
choice of \(|W|=\lceil3\Gamma^2x^{C_H}\rceil\) distinct bit strings
of length \(k-2\)
\releasecite{Development/AbelianVarietyTheory/Proofs/GlobalInterpolation/MPowFiveMulFMulLLe.lean}{40}{Injective words}.
Such a choice exists once \(3\Gamma^2x^{C_H}+1\le2^{k-2}\). With
\(\Gamma=\exp(C_\Gamma m^5f\log b)\), Section~\ref{sec:family} gives
\(m^5f\log b\le128kx^{-(21/16)\varepsilon}\). Hence
\[
 \log(3\Gamma^2x^{C_H}+1)\le2+256C_\Gamma kx^{-(21/16)\varepsilon}
   +C_H\log x\le(k-2)\log2
\]
beyond a threshold that depends only on \(C_\Gamma,C_H,\varepsilon\),
because \(k\ge x^{9/2+10\varepsilon}\). The words are inputs of the
records in (e), so every index has
\(3\Gamma^2x^{C_H}\le|W|\le3\Gamma^2x^{C_H}+1\).

\paragraph{What the selection takes from the release.}
\teamanchor{proof.selection-residual}%
The edition takes the following from the release.
\begin{enumerate}
\item[(R1)] \emph{Record existence}
\releasecite{Development/AbelianVarietyTheory/Definitions/GlobalTwisting/ValidForConstructedPointClassesEventuallyWorking.lean}{190}{Record existence}.
For the precision and the constants fixed in (a), taken large
enough, there are \(c_s>0\) and \(x_0\) with the following property. For every scale \(x\ge x_0\) and every ranks, host, arrays
and words satisfying (b)--(d), there are column data
\(d_w^{(0)},D_\star>0\) with \(\log Y_w\ge c_smx^\varepsilon\) for every
word. Moreover, for every choice of primes \(T_w\) in the windows with
norm generators \(\equiv1\bmod Q_Sqd(g_0g_1)^2\), with
\(T_w\equiv1\bmod Q_Sqd\) and \(T_w\ne p,p'\), there is a record
over these levels. It satisfies the record, ledger, local-twist and
validity conditions, and it carries the class property: the rows of
item~2, the field \(K_0\), its admissibility at both auxiliary levels
and \(s(K_0)\le P_O\). The edition derives the bound on \(\log Y_w\)
(with \(c_s=14\)), the bounds \(b^{m/2}\le T_w\) and
\(\log T_w\le C_Tm\log b\) and the bounds on \(D_\Gamma\) above, once
the column data are the mass--Vandermonde products of
Section~\ref{sec:integral}. It takes from the release the existence of
the record over the chosen levels, that identification of the column
data, and the class property.
\item[(R2)] In the release, (R1) is proved from the class property by
\releasecite{Development/AbelianVarietyTheory/Proofs/IwasawaDecomposition/MemberTopFinalContextParameters.lean}{14}{Record existence from the context parameters}.
That proof turns volumes into masses
\releasecite{Development/AbelianVarietyTheory/Proofs/IwasawaDecomposition/VolumesEqMassesOfUnfold.lean}{12}{Volumes equal masses}
and uses the following inputs.
\begin{itemize}
\item The mass ratio for unramified CM fields
\releasecite{Development/AbelianVarietyTheory/Proofs/GlobalInterpolation/MassRatioAllUnramifiedCmFields.lean}{8}{Mass ratio}.
For the standard self-dual genus this is the successive mass quotient
\eqref{eq:arithmetic-mass-quotient}, which Section~\ref{sec:integral}
derives from the exact mass formula of Gan--Hanke--Yu and the
functional equation.
\item A hyperplane volume ratio. Let \(W\) be a nondegenerate
hermitian space over \(E/F\) of rank \(M\ge3\) with the same
signature at every real place, and let \(\pm\) be a sign that occurs in
that signature. There is \(x\in W\) with \(h(x,x)=\pm1\) such that
\(W'=x^\perp\) satisfies, in the volume normalization of the release and under its
volume tie,
\[
 \frac{\operatorname{vol}(W')}{\operatorname{vol}(W)}
  =\left(\frac{(2\pi)^M}{(M-1)!}\right)^{\!f}
   (2v_0)^{f(1/2-M)}\,L_F(M,\eta_{E/F}^{\,M})^{-1}
\]
\releasecite{Formalism/SharpRestriction/GlobalSchurBounds/ThickPlacesCorollaryOfAnyTranslate.lean}{2330}{The hyperplane volume ratio}.
The class selection receives it as a hypothesis
\releasecite{Development/AutomorphicArithmetic/ThetaCorrespondence/ThetaPairLifts/ThirdPlaceSeries/NonemptyHypsOfAllRows.lean}{46}{Volume ratio passed to the class selection}.
The hypothesis is discharged at
\releasecite{Development/AutomorphicArithmetic/ThetaCorrespondence/ThetaPairLifts/ThirdPlaceSeries/HypothesesNonemptyOfUnramifiedExponents.lean}{60}{Discharge of the volume ratio}.
There the release unfolds the volume of the orbit of \(x\) through the
one-vector Siegel--Weil identity
\releasecite{Development/AbelianVarietyTheory/Proofs/FlagMeasures/UnitaryHyperplaneVolumeRatioSiegelWeil.lean}{10}{Volume ratio from Siegel--Weil}
\releasecite{Development/AbelianVarietyTheory/Proofs/GlobalInterpolation/VolumeRatio.lean}{14}{Unfolding and the local densities}.
The identity used is the Siegel--Weil formula in the convergent range
(finite volume and equality of the theta integral with the product of
local densities). The release proves it
\releasecite{CompanionTheory/SpectralGeometry/LocalGlobalProofs/AdelicAnalysis/Imports.lean}{10}{Convergent Siegel--Weil formula}
and passes it through its literature bundle
\releasecite{Development/AutomorphicArithmetic/ThetaCorrespondence/ThetaPairLifts/ThirdPlaceSeries/EveryCharacterOfLiteratureOfVWin.lean}{708}{Supply of the convergent Siegel--Weil formula}.
For unitary groups compare \cite[Theorem~2 and the convergence
statements preceding it]{tr3yamana}; the edition does not match the
normalizations here. The finite factors can be checked
directly. At a finite place \(v\) of \(F\), with residue
cardinality \(q\), the lattices are self-dual and \(E/F\) is unramified
or split. A unit-norm vector modulo the uniformizer lifts by Hensel's
lemma, since the reduced form is nondegenerate. In the inert case a
nondegenerate hermitian space of dimension \(M\) over
\(\mathbf F_{q^2}\) has \(q^{M-1}(q^M-(-1)^M)\) vectors of a given
nonzero norm. In the split case the equation \(x\cdot y=a\) on
\(\mathbf F_q^M\times\mathbf F_q^M\) has \(q^{M-1}(q^M-1)\) solutions.
Dividing by \(q^{2M-1}\) gives the local densities \(1-(-q)^{-M}\) and
\(1-q^{-M}\). Both equal \(L_v(M,\eta_{E/F}^{\,M})^{-1}\), and their
product is the \(L\)-value above. The real factors and the
discriminant power \((2v_0)^{f(1/2-M)}\) come from the measure
normalizations of the release and are not rederived here.
\item The group-retwisting conditions and the bundle conditions
\(\mathrm{ExBL}\) and \(\mathrm{RowL}\). These are supplied from the
release's third-place hypothesis bundle at
\releasecite{Development/AutomorphicArithmetic/ThetaCorrespondence/ThetaPairLifts/ThirdPlaceSeries/HypothesesNonemptyOfUnramifiedExponents.lean}{41}{Supply of the selection hypotheses}.
This edition identifies them only where they are used elsewhere in the
article.
\end{itemize}
\item[(R3)] The arrays of (c). Their windows are nonempty by (c). The
release chooses the arrays inside these windows so that they also
satisfy its admissibility condition
\[
 \sum_{p<X',\ p\ \text{candidate}}\frac{\log p}{p}
   \Bigl(2\cdot\mathbf 1[\,x^{-\varepsilon}<\mathrm{Sig}_p\,]+4\,\mathrm{Sig}_p\Bigr)
   \le x^{-9}
\]
\releasecite{Development/AbelianVarietyTheory/Definitions/GlobalTwisting/ThinScalesIffHostBound.lean}{149}{The admissibility condition}.
Here the candidates and the signature excess \(\mathrm{Sig}_p\) are
those of the local prime tables of the release. Such arrays exist at
every point of the family by a mean-excess argument
\releasecite{Development/AbelianVarietyTheory/Proofs/GlobalInterpolation/WindowShareRequirementClosedFormProof.lean}{17}{Weight vectors at every point}.
They are transferred to the host
\releasecite{Development/AbelianVarietyTheory/Proofs/GlobalInterpolation/WeightWindowRequirementHostOfPoint.lean}{17}{Arrays on the host}
and attached to it
\releasecite{Development/AutomorphicArithmetic/CharacterBoxes/HostWW.lean}{27}{Hosts with admissible arrays}.
This edition does not rederive that argument. Also the host data other
than \(g_0,g_1,p,p'\). The latter are hypotheses of the statement,
from Section~\ref{sec:family}.
\end{enumerate}
Everything else in items~1--3 follows from (R1)--(R3) and the
derivations above.

\paragraph{The rows and field on which the two section estimates hold.}
\teamanchor{proof.section-run-family-identification}%
Section~\ref{sec:integral} constructs the two auxiliary section
systems on the base level \(g_0g_1\) with their respective additional
primes. Its numerical estimates use a common geometric coefficient
field \(K_0\). This field is chosen before the larger field containing
the moving character values. They also use the Harish--Chandra rows
with specified spacings. Consecutive entries in the initial block
differ by \(4\Delta_{\rm cl}\). In each subsequent block the nodes are
\(B_\tau+lH_{\rm node}\) and \(B_\tau+(l+\tfrac12)H_{\rm node}\), with
the same prescribed \(\Delta_{\rm cl},H_{\rm node}\) for all words. These choices give the gaps in the flag
weights and the bound on the degree and root discriminant of \(K_0\)
used in the section heights.

The identification is simultaneous, and the release supplies it. The
selection is assembled from points constructed over contexts. A
context fixes \(\varepsilon\), \(u_{\rm cl}\), \(u_{\rm node}\), the host, the
words and the rows above before the base level is built, and it
requires those row values for every word
\releasecite{Development/AbelianVarietyTheory/Definitions/GlobalTwisting/SetUnipotentTranslate.lean}{37}{Rows fixed in the context}.
A point over the context therefore has stage parameter \(\Delta_{\rm cl}\)
\releasecite{Development/AbelianVarietyTheory/Definitions/AutomorphicMembership/ExistsStagesClauses.lean}{94}{The stage parameter}.
At the base level the column choice fixes field data and requires the
base coefficient field to equal the field they generate. This happens
before the later stages are built. That field is \(K_0\). It is
admissible at both auxiliary levels and contains the frame character
values
\releasecite{Development/AbelianVarietyTheory/Proofs/IwasawaDecomposition/IndexBoundCoefficientFieldOfExactSmallField.lean}{28}{Admissibility of the exact base field}.
It satisfies \(s(K_0)\le P_O\) beyond one threshold \(x_0\), which
depends only on \(\varepsilon\), the family constants and
\(u_{\rm cl},u_{\rm node}\), and is fixed before the scale
\releasecite{Development/AbelianVarietyTheory/Proofs/HeightProofs/EventuallyConstructedPointHostSizeNumerical.lean}{19}{Eventual size of the field}.

Each family index beyond the family threshold is such a constructed
point
\releasecite{Development/AbelianVarietyTheory/Definitions/GlobalTwisting/ValidForConstructedPointClassesEventuallyWorking.lean}{250}{Indices are constructed points},
and one choice of the context exponents serves the whole family
\releasecite{Development/AbelianVarietyTheory/Definitions/GlobalTwisting/ValidForConstructedPointClassesEventuallyWorking.lean}{264}{One selection for every index}
\releasecite{Development/AutomorphicArithmetic/CharacterBoxes/ClassSel.lean}{14}{Its use in the class selection}.
Each such index therefore carries the rows of its point
\releasecite{Development/AbelianVarietyTheory/Definitions/GlobalTwisting/PairRowIffQdRows.lean}{107}{The row values carried by a point},
and it inherits the two field admissibilities and the bound on
\(s(K_0)\) through the ties established for every bundle
\releasecite{Development/AutomorphicArithmetic/ThetaCorrespondence/ThetaPairLifts/ThirdPlaceSeries/LinkedTies.lean}{32}{Ties at every bundle}.
The periods, test columns and tied records at an index belong to the
same point. Let \(i_0\) be a seed and \(u\in H\). The index \(i(u)\)
of item~3 is a family index. It carries its own rows, its own \(K_0\)
at both auxiliary levels and its own periods and test columns, and
the values \(\Delta_{\rm cl},H_{\rm node}\) are common to all of them. The two runs are
built at that point from these data
\releasecite{Development/AutomorphicArithmetic/ThetaCorrespondence/ThetaPairLifts/ThirdPlaceSeries/ExistsRunsAt.lean}{18}{The two runs at one point}.
The section estimates of Section~\ref{sec:integral} therefore apply at
every selected index. The identification is supplied by the release
and has been traced there. In the order of construction above it is
part of the record existence (R1), and it has not been rederived in
this edition.

\paragraph{The integral lattice in the ramified columns.}
\teamanchor{proof.ramified-donor-identification}%
Fix a word \(w\) at its split level prime \(T_w\). In the complex
column space, Section~\ref{sec:trace} computes the subspace on which
the averages in all \(f\) residue-unit directions vanish. Let \(V_r\)
be its rational form in \(M_w\otimes K\). Its dimension is
\[
 d_r=d_M\left(\frac{T_w-2}{T_w-1}\right)^f,\qquad
 d_M=\operatorname{rk}M_w .
\]
Put \(H_w=M_w\cap V_r\), with the metrics of \(M_w\) restricted to
\(V_r\) at every complex embedding
\releasecite{Formalism/LocalRingCardinality/DataForms.lean}{243}{The sum-zero lattice}.
It is saturated in \(M_w\). On this sum-zero model the Gram
determinant gives \(\mu(H_w)-\mu(M_w)=-f\log(T_w-1)/(T_w-2)\), as
recorded in
\releasecite{Formalism/SharpRestriction/Divisibility/LocalQuadraticParts.lean}{621}{The ramified rank and metrics}.

The undivided source \(N_w\) is the sublattice of \(H_w\) determined
by local conditions. At a served place \(v\nmid T_w\), \((N_w)_v\) is
the preimage of the local period lattice \(N_{P,v}\). At \(v\mid T_w\)
it is the preimage of \(T_w^{e_T}U^m(S_{0,w}\cap V_r)\). At every other
place it is \((H_w)_v\). The release carries the first two
characterizations with the selected data, as fields of every index
beyond the family threshold
\releasecite{Formalism/LocalRingCardinality/DataLatDefs.lean}{42}{The condition at a served place}
\releasecite{Formalism/LocalRingCardinality/DataLatDefs.lean}{111}{The condition at the level prime}
\releasecite{Formalism/LocalRingCardinality/DataLatDefs.lean}{115}{The local characterizations together}
\releasecite{Formalism/LocalRingCardinality/DataExtra.lean}{51}{Their place in the index data}
\releasecite{Formalism/SharpRestriction/ParConsts.lean}{77}{Every index carries them}.
The third follows from the unserved case below. A full-rank lattice is
determined by its localizations. Its rational space
is \(V_r\), and every column record identifies the source space with
\(V_r\) isometrically for the zero-stage norm
\releasecite{Formalism/LocalRingCardinality/RhoDetPrimeEqMuTildeMulChiPow.lean}{291}{The zero-stage isometry of the source}
\releasecite{Development/AbelianVarietyTheory/Definitions/AutomorphicMembership/ExistsDivStage.lean}{334}{The source space and its embedding}.
Hence \(N_w\subseteq H_w\) is a full-rank isometric inclusion. The
source is not spanned by chosen test vectors. Section~\ref{sec:periods}
starts its descent from an arbitrary element of this lattice, so no
further comparison with another lattice is needed.

Let \(v\) be a coefficient place of absolute ramification index
\(e_v\). If \(v\nmid T_w\) is unserved, then \(\tau_v=\iota_v=0\) and
\((N_w)_v=(H_w)_v\)
\releasecite{Formalism/SharpRestriction/GrowthBounds/PlaneCountTheoremWithConstants.lean}{720}{The unserved places}.
If \(v\nmid T_w\) is served, then \(\tau_v\) is a
nonnegative integer and the source satisfies
\[
 (N_w)_v\subseteq\mathfrak p_v^{e_v\tau_v}(H_w)_v,\qquad
 \operatorname{length}_{\mathcal O_v}((H_w)_v/(N_w)_v)
     \le e_v d_M(\tau_v+\iota_v).
\]
The \hyperlink{proof.served-length-bound}{next paragraph} derives
both bounds from the release's lattice statement at the individual
sample vertices. The gain valuation on the source at a served place
is taken from the release
\releasecite{Formalism/SharpRestriction/GrowthBounds/SmallOrderPrimes.lean}{220}{Inclusion, length and gain at a served place}
\releasecite{Formalism/SharpRestriction/GrowthBounds/TopChoiceBounds.lean}{1586}{The gain on the local period lattice}.
At \(v\mid T_w\) we put \(\tau_v=0\). There \((N_w)_v\subseteq(H_w)_v\), and
the lengths satisfy the weighted level-index bound used in the
\hyperlink{proof.divided-column-slopes}{calculation of the
divided-source slopes}
\releasecite{Formalism/SharpRestriction/Divisibility/LocalQuadraticParts.lean}{621}{The ramified rank, metrics and level index}.
The \hyperlink{proof.level-index-bound}{level-index paragraph} below
derives this bound from the release's local Hecke relations at the
vertices over \(T_w\) and its size data.
These local facts hold for every family in the selected class. Its
twist selection and the duty conditions of its members are part of
the
\releasecite{Development/AutomorphicArithmetic/CharacterBoxes/ClassSel.lean}{14}{class selection theorem}.
The release assembles them in
\releasecite{Development/AutomorphicArithmetic/ThetaCorrespondence/ThetaPairLifts/ThirdPlaceSeries/LinkG.lean}{13}{the closed comparison for the class},
which lies on the path to the final theorem
\releasecite{Development/AutomorphicArithmetic/ThetaCorrespondence/ThetaPairLifts/ThirdPlaceSeries/TopShape.lean}{54}{Its use in the top shape}.

\paragraph{The inclusion and length bound at a served place.}
\teamanchor{proof.served-length-bound}%
Fix a served place \(v\nmid T_w\) over \(p\). Let \(\mathcal O\) be the
discrete valuation ring of the local coefficient field in which the
release builds its local period model, and put \(n=v_{\mathcal O}(p)\),
so that \(\operatorname{length}_{\mathcal O}(\mathcal O/p\mathcal O)=n\).
Then \(n=e_ve'\), where \(e'\) is the ramification index of
\(\mathcal O\) over \(\mathcal O_v\). The comparison map \(c\) from
\(M_w\otimes K\) to the local column space is semilinear over an
embedding of \(K\) with \(K\cap\mathcal O=\mathcal O_{(v)}\), and it carries a
basis of \(M_w\otimes K\) to a linearly independent family. The
release's local lattice \(M\) is the \(\mathcal O\)-span of \(c(M_w)\)
\releasecite{Formalism/LocalRingCardinality/DataLatDefs.lean}{26}{The local lattice is spanned by the source}.
It is free of rank \(d_M\). Write \(V_r\) also for the span of
\(c(V_r)\), and put \(H=M\cap V_r\). It is saturated in \(M\) because
\(V_r\) is a subspace. Flat base change identifies \(H\) with the
\(\mathcal O\)-span of \(c(H_w)\)
\releasecite{Formalism/LocalRingCardinality/Lemmas.lean}{968}{The ramified part of the local model}.

The release attaches to \(v\) the following data
\releasecite{Formalism/SharpRestriction/GrowthBounds/IsResidueClassZero.lean}{292}{The served samples}
\releasecite{Formalism/SharpRestriction/GrowthBounds/IsResidueClassZero.lean}{341}{The restriction loss}
\releasecite{Formalism/SharpRestriction/GrowthBounds/IsResidueClassZero.lean}{348}{The index allowance}.
There is a nonempty finite set \(V_v\) of good sample orbits. Each
\(s\in V_v\) has a vertex operator \(y_s\), and these operators are
commuting \(\mathcal O\)-linear endomorphisms. There are integers
\(t_s,D_v,S_s,\varepsilon_v\ge0\) and real weights \(\kappa_s\ge0\).
At a served place
\(\Sigma_v=m\sum_{s\in V_v}\kappa_s\le x^{-\varepsilon}\), and \(m\ge1\). Hence
\(\Sigma_v\le\frac12\) once \(x^{-\varepsilon}\le\frac12\). With \(S=\max_sS_s\), the
loss and the allowance are
\[
 \tau_v=\varepsilon_v+\tau_v',\qquad
 \tau_v'=\sum_{s\in V_v}m(t_s+D_v+1),\qquad
 \iota_v=2S\Sigma_v .
\]
Thus \(\tau_v\) is a nonnegative integer. For \(B\subseteq V_v\), put
\(Y_B=\prod_{s\in B}y_s^m\) and \(\sigma(B)=\sum_{s\in V_v\setminus B}S_s\).
For an \(\mathcal O\)-lattice \(L\), let
\(N(L)=\sum_{B\subseteq V_v}p^{\sigma(B)}Y_BL\). The local period lattice is
\[
 N_{P,v}=N(M)\cap p^{\tau_v}M\cap V_r
\]
\releasecite{Formalism/SharpRestriction/GrowthBounds/IsResidueClassZero.lean}{464}{The local period lattice}.

Two release statements enter. First, there is a lattice
\(M^{\rm na}\) with \(p^{\varepsilon_v}M^{\rm na}\subseteq M\subseteq
M^{\rm na}\), and the \(y_s\) preserve both lattices
\releasecite{Formalism/SharpRestriction/GrowthBounds/CornerGLContinuity.lean}{191}{The lattice sandwich}
\releasecite{Formalism/SharpRestriction/GrowthBounds/PlaneCountTheoremWithConstants.lean}{1529}{The operators preserve the lattices}.
The sandwich forces \(\operatorname{rk}M^{\rm na}=d_M\). Second, the
good-rank property holds at every served vertex with the parameters
\(t_s,D_v,\kappa_s\)
\releasecite{Formalism/SharpRestriction/GrowthBounds/CornerGLContinuity.lean}{180}{The good-rank property at the served vertices}.
The release deduces from it, together with the trivial stabilizers
and finite classes of the local lattice tuples, an all-but statement. For every
\(\beta\ge0\) there is a finite set \(C_\beta\subseteq M^{\rm na}\) with
\[
 |C_\beta|\le d_Mm^{\beta+1}e_{\beta+1}(\kappa),\qquad
 p^{\tau_v'+\beta S}M^{\rm na}\subseteq N(M^{\rm na})
   +\mathcal O C_\beta ,
\]
where \(e_k\) is the elementary symmetric function of the weights
\((\kappa_s)_{s\in V_v}\)
\releasecite{Formalism/SharpRestriction/GrowthBounds/SelectionLemmas.lean}{41}{The all-but statement at a vertex}
\releasecite{Formalism/SharpRestriction/GrowthBounds/LatticeSystems.lean}{4417}{Its proof from the good-rank property}.
The \hyperlink{proof.all-but-walk}{walk argument} after this
paragraph derives the statement from the release's local lattice
system at the vertices. The rest of the argument is derived here.

The operators \(Y_B\) are \(\mathcal O\)-linear, so
\(p^{\varepsilon_v}p^{\sigma(B)}Y_Bz=p^{\sigma(B)}Y_B(p^{\varepsilon_v}z)\),
and \(p^{\varepsilon_v}z\in M\) for \(z\in M^{\rm na}\). Hence
\(p^{\varepsilon_v}N(M^{\rm na})\subseteq N(M)\). Multiply the all-but
inclusion by \(p^{\varepsilon_v}\) and use \(M\subseteq M^{\rm na}\).
With \(C'_\beta=p^{\varepsilon_v}C_\beta\subseteq M\), this gives
\[
 p^{\tau_v+\beta S}M\subseteq N(M)+\mathcal O C'_\beta
 \qquad(\beta\ge0).
\]
Also \(N(M)\subseteq M\), and the term \(B=\varnothing\) gives
\(p^{\sigma(\varnothing)}M\subseteq N(M)\). Since the weights are
nonnegative, every monomial of \(e_k(\kappa)\) occurs \(k!\) times in
\((\sum_s\kappa_s)^k\). Hence
\(m^ke_k(\kappa)\le\Sigma_v^k/k!\) and
\(|C_\beta|\le d_M\Sigma_v^{\beta+1}/(\beta+1)!\).

Put \(Q=M/N(M)\). It is killed by \(p^{\sigma(\varnothing)}\), so it has
finite length. Let \(Q_\beta\) be the image of \(p^{\tau_v+\beta S}M\). If
\(S=0\), then \(\sigma(\varnothing)=0\) and \(Q=0\). Otherwise
\(Q_{\beta+1}=p^SQ_\beta\), and \(Q_\beta=0\) once
\(\beta S\ge\sigma(\varnothing)\). The module \(Q_\beta\) lies in the
image of \(\mathcal OC'_\beta\). Over the discrete valuation ring
\(\mathcal O\), a submodule of a module with \(r\) generators has at most
\(r\) generators. Therefore
\[
 \operatorname{length}_{\mathcal O}(Q_\beta/p^SQ_\beta)\le |C_\beta|\,nS,
 \qquad
 \operatorname{length}_{\mathcal O}Q_0
   \le nSd_M\sum_{k\ge1}\frac{\Sigma_v^k}{k!}
   =nSd_M(e^{\Sigma_v}-1)\le 2nSd_M\Sigma_v ,
\]
using \(e^y-1\le2y\) for \(0\le y\le\frac12\). Now
\(p^{\tau_v}M/(N(M)\cap p^{\tau_v}M)\cong Q_0\) and
\(\operatorname{length}_{\mathcal O}(M/p^{\tau_v}M)=nd_M\tau_v\). Hence
\[
 \operatorname{length}_{\mathcal O}\bigl(M/(N(M)\cap p^{\tau_v}M)\bigr)
   \le nd_M(\tau_v+\iota_v).
\]

Now pass to the ramified part. If \(x\in N_{P,v}\), then
\(x=p^{\tau_v}y\) with \(y\in M\), and \(y\in V_r\). Thus \(y\in H\) and
\(N_{P,v}\subseteq p^{\tau_v}H\). Moreover
\(N_{P,v}=H\cap(N(M)\cap p^{\tau_v}M)\), so \(H/N_{P,v}\) embeds in
\(M/(N(M)\cap p^{\tau_v}M)\). Hence
\(\operatorname{length}_{\mathcal O}(H/N_{P,v})\le nd_M(\tau_v+\iota_v)\).

It remains to descend to the source. By definition
\((N_w)_{(v)}\) consists of the \(x\in(H_w)_{(v)}\) with \(c(x)\in
N_{P,v}\)
\releasecite{Formalism/LocalRingCardinality/DataLatDefs.lean}{42}{The condition at a served place}.
For the inclusion, let \(x\in(N_w)_{(v)}\). Then
\(c(p^{-\tau_v}x)\in H\subseteq M\). The coordinates of
\(p^{-\tau_v}x\) in an \(\mathcal O_{(v)}\)-basis of \((M_w)_{(v)}\) lie
in \(K\). Their images lie in \(\mathcal O\), because \(c\) of this basis
is linearly independent and spans \(M\). They
therefore lie in \(\mathcal O_{(v)}\). Since
\(p^{-\tau_v}x\in V_r\), we get
\(p^{-\tau_v}x\in(M_w)_{(v)}\cap V_r=(H_w)_{(v)}\). Because
\(p\mathcal O_v=\mathfrak p_v^{e_v}\), this proves
\((N_w)_v\subseteq\mathfrak p_v^{e_v\tau_v}(H_w)_v\).

For the length bound one needs
\(N_{P,v}=\mathcal O\,c((N_w)_{(v)})\). The inclusion \(\supseteq\) is
immediate from the definition. The release proves the reverse
inclusion. It shows that \(N_{P,v}\) is spanned by \(c\) of its
\(K\)-rational elements, because each \(y_s\) is induced by a
\(K\)-linear endomorphism of the rational column space
\releasecite{Formalism/LocalToGlobal/GroupLattices/ComparisonMaps/YvRatOfValRat.lean}{199}{The period coordinates are rational}
\releasecite{Formalism/SharpRestriction/GrowthBounds/JointsLit.lean}{256}{The period lattice is spanned by rational elements}
\releasecite{Formalism/LocalToGlobal/GroupLattices/ComparisonMaps/NpDescOfValRat.lean}{13}{Its use for the source}.
Granting this, \(\mathcal O\) is flat over \(\mathcal O_{(v)}\), and
\[
 \mathcal O\otimes_{\mathcal O_{(v)}}\bigl((H_w)_{(v)}/(N_w)_{(v)}\bigr)
   \cong H/N_{P,v}.
\]
For a module \(Q'\) of finite length over \(\mathcal O_{(v)}\), the
\(\mathcal O\)-length of \(\mathcal O\otimes Q'\) is \(e'\) times the
length of \(Q'\), since \(\mathcal O/\mathfrak p_v\mathcal O\) has
length \(e'\). Dividing by \(e'\), we get
\[
 \operatorname{length}_{\mathcal O_v}((H_w)_v/(N_w)_v)
   \le e_vd_M(\tau_v+\iota_v).
\]
The release records the same two bounds for the source in
\releasecite{Formalism/SharpRestriction/GrowthBounds/PlaneCountTheoremWithConstants.lean}{664}{The inclusion for the source}
\releasecite{Formalism/SharpRestriction/GrowthBounds/PlaneCountTheoremWithConstants.lean}{720}{The length for the source}.

\paragraph{The all-but statement from the local lattice system.}
\teamanchor{proof.all-but-walk}%
Keep the served place \(v\) and its data. The release's local ring
\(\mathcal O\) contains an element \(\varpi\) with \(\varpi^2=p\); this is a field of
the release's local setting at every vertex
\releasecite{Formalism/SharpRestriction/GrowthBounds/UnramifiedQuadraticSetting.lean}{29}{The square root of \(p\) in the local ring}. A
labelling is a map \(\gamma\) from \(V_v\) to the three labels
base, ball and dual. For \(s\in V_v\) and a label \(\ell\), write
\(\gamma[s\mapsto\ell]\) for the labelling that agrees with \(\gamma\)
off \(s\) and has \(\ell\) at \(s\). For \(B\subseteq V_v\), let
\(\gamma_B\) be the labelling with dual on \(B\) and base elsewhere.
The release attaches to every labelling an \(\mathcal O\)-lattice
\(L_\gamma\) in a common \(\mathcal O\)-module on which the \(y_s\)
act. For all \(\gamma\) and \(s\in V_v\) these lattices satisfy:
\begin{enumerate}
\item[(i)] \(L_{\gamma_\varnothing}=M^{\rm na}\);
\item[(ii)] \(L_{\gamma[s\mapsto{\rm base}]}\subseteq L_{\gamma[s\mapsto{\rm ball}]}\) and
 \(L_{\gamma[s\mapsto{\rm dual}]}\subseteq L_{\gamma[s\mapsto{\rm base}]}\);
\item[(iii)] \(\varpi^{t_s}L_{\gamma[s\mapsto{\rm ball}]}\subseteq L_{\gamma[s\mapsto{\rm base}]}\) and
 \(\varpi^{t_s}L_{\gamma[s\mapsto{\rm base}]}\subseteq L_{\gamma[s\mapsto{\rm dual}]}\);
\item[(iv)] for every nonempty \(B\subseteq V_v\) there is a finite set
 \(C_B\subseteq L_{\gamma_B}\) with \(|C_B|\le d_M\prod_{s\in B}\kappa_s\) and
 \[
  p^{D_v+1}L_{\gamma_B}\subseteq\sum_{s\in B}y_s\,L_{\gamma_B[s\mapsto{\rm ball}]}
   +\mathcal O C_B .
 \]
\end{enumerate}
\releasecite{Formalism/SharpRestriction/GrowthBounds/LatticeSystems.lean}{197}{The lattice system at the vertices}
\releasecite{Formalism/SharpRestriction/GrowthBounds/LatticeSystems.lean}{2465}{The local chain of lattices}
\releasecite{Formalism/SharpRestriction/GrowthBounds/LatticeSystems.lean}{4300}{The generation statement from the orbit decomposition}
\releasecite{Formalism/SharpRestriction/GrowthBounds/LatticeSystems.lean}{3746}{The local Smith form from the good-rank property}
\releasecite{Formalism/SharpRestriction/GrowthBounds/LatticeSystems.lean}{4194}{The rank count}
\releasecite{Formalism/SharpRestriction/GrowthBounds/LatticeSystems.lean}{2886}{Trivial stabilizers and finite classes give the decomposition}.
The lattices consist of functions on the local lattice tuples. The
generation statement~(iv) is proved in the release by splitting a
function into its orbits under the level groups, which uses the trivial
stabilizers and the finite classes, by a Smith-form computation at each
vertex of \(B\) from the good-rank property, and by counting the orbit
representatives against \(\operatorname{rk}M^{\rm na}=d_M\). The edition
takes (i)--(iv) from the release. The all-but statement follows from
them by the walk argument below, which is the release's abstract
counting lemma
\releasecite{Formalism/SubspaceCounting/Statement.lean}{38}{The walk statement}
\releasecite{Formalism/SubspaceCounting/SubspaceCountingProofs/WalkSharp.lean}{722}{Its proof}.

The \(y_s\) commute, are \(\mathcal O\)-linear and preserve \(M^{\rm na}\).
For \(e\in\{0,\dots,m\}^{V_v}\) put \(y^e=\prod_sy_s^{e_s}\). If
\(E=\{s:e_s=m\}\), then \(y^eM^{\rm na}\subseteq Y_EM^{\rm na}\), and therefore
\[
 p^{\sigma(E)}y^eM^{\rm na}\subseteq N(M^{\rm na}).
 \tag{$\ast$}
\]
Put \(V=|V_v|\). If \(\beta+1>V\), then
\(\sigma(\varnothing)\le VS\le\beta S\), so
\(p^{\tau_v'+\beta S}M^{\rm na}\subseteq N(M^{\rm na})\) and
\(C_\beta=\varnothing\) works. Let \(\beta+1\le V\), and number the
vertices \(s_0,\dots,s_{V-1}\). A state is a pair \((B,i)\) with
\(B\subseteq V_v\), \(|B|=\beta+1\), and \(i:B\to\{0,\dots,m-1\}\). Let
\(\operatorname{top}(B)\) be the largest index of a vertex in \(B\). The
exponent \(e=e(B,i)\) of the state has \(e_s=i_s\) on \(B\), \(e_s=m\) for
\(s\notin B\) of index below \(\operatorname{top}(B)\), and \(e_s=0\) for
\(s\) of index above \(\operatorname{top}(B)\). Its budget is
\[
 R(B,i)=\sum_{s\in B}\bigl(t_s+2(m-1-i_s)t_s+2(m-i_s)(D_v+1)\bigr)
   +\sum_{s\text{ above }\operatorname{top}(B)}2m(t_s+D_v+1).
\]
Let \(C_\beta\) be the set of all \(y^{e(B,i)}c\) with \((B,i)\) a state
and \(c\in C_B\). By (i) and (ii), \(C_B\subseteq L_{\gamma_B}\subseteq
M^{\rm na}\), so \(C_\beta\subseteq M^{\rm na}\). There are
\(m^{\beta+1}\) states for each \(B\), hence
\(|C_\beta|\le\sum_{|B|=\beta+1}m^{\beta+1}d_M\prod_{s\in B}\kappa_s
=d_Mm^{\beta+1}e_{\beta+1}(\kappa)\).

We show that for every state \((B,i)\) and every \(z\in L_{\gamma_B}\),
\[
 \varpi^{R(B,i)}p^{\beta S}y^{e(B,i)}z\in N(M^{\rm na})+\mathcal OC_\beta .
\]
The induction runs downward on \(\sum_se(B,i)_s\le mV\). Write
\(R=R(B,i)\) and \(e=e(B,i)\). Since \(B\) is nonempty and \(i_s<m\), we
have \(R\ge2(D_v+1)\). By (iv),
\(p^{D_v+1}z=\sum_{s\in B}y_sz_s+c\) with
\(z_s\in L_{\gamma_B[s\mapsto{\rm ball}]}\) and \(c\in\mathcal OC_B\). Hence
\(\varpi^Rp^{\beta S}y^ez=\varpi^{R-2(D_v+1)}p^{\beta S}\bigl(\sum_sy^{e+1_s}z_s+y^ec\bigr)\),
where \(1_s\) is the indicator of \(s\). The term \(y^ec\) lies in
\(\mathcal OC_\beta\). For the term of \(s\in B\) there are three cases.
\begin{enumerate}
\item[(a)] \(i_s+1<m\). By (iii) twice, \(\varpi^{2t_s}z_s\in L_{\gamma_B}\).
 The state \((B,i+1_s)\) has exponent \(e+1_s\) and budget
 \(R-2(D_v+1)-2t_s\). The term equals
 \(\varpi^{R(B,i+1_s)}p^{\beta S}y^{e+1_s}(\varpi^{2t_s}z_s)\), which lies in
 \(N(M^{\rm na})+\mathcal OC_\beta\) by induction.
\item[(b)] \(i_s+1=m\) and \(\operatorname{top}(B)<V-1\). Let \(u\) be the
 vertex following \(\operatorname{top}(B)\), and put
 \(B'=(B\setminus\{s\})\cup\{u\}\), with \(i'=i\) on \(B\setminus\{s\}\) and
 \(i'_u=0\). By (iii), \(\varpi^{t_s}z_s\in L_{\gamma_{B\setminus\{s\}}}\) and
 \(\varpi^{t_s+t_u}z_s\in L_{\gamma_{B'}}\). Now \(s\) lies below
 \(u=\operatorname{top}(B')\), so \(e(B',i')=e+1_s\). The budget is
 \(R(B',i')=R-2(D_v+1)-t_s-t_u\). The term equals
 \(\varpi^{R(B',i')}p^{\beta S}y^{e(B',i')}(\varpi^{t_s+t_u}z_s)\) and is
 covered by induction.
\item[(c)] \(i_s+1=m\) and \(\operatorname{top}(B)=V-1\). By (iii) and (ii),
 \(\varpi^{t_s}z_s\in L_{\gamma_{B\setminus\{s\}}}\subseteq M^{\rm na}\). The
 exponent \(e+1_s\) equals \(m\) off \(B\setminus\{s\}\), a set of
 \(\beta\) vertices. Hence \(\sigma(E)\le\beta S\) for its set \(E\), and
 \(p^{\beta S}y^{e+1_s}(\varpi^{t_s}z_s)\in N(M^{\rm na})\) by
 \((\ast)\). The remaining factor \(\varpi^{R-2(D_v+1)-t_s}\) has a
 nonnegative exponent.
\end{enumerate}
In (a) and (b) the exponent sum grows by one, so the induction is
well founded. Finally let \(x\in M^{\rm na}\), and take
\(B_0=\{s_0,\dots,s_\beta\}\) with \(i=0\). By (i) and (iii),
\(\varpi^{\sum_{s\in B_0}t_s}x\in L_{\gamma_{B_0}}\), and
\(e(B_0,0)=0\). The budget is
\(R(B_0,0)=2\tau_v'-\sum_{s\in B_0}t_s\). Hence
\[
 p^{\tau_v'+\beta S}x
 =\varpi^{R(B_0,0)}p^{\beta S}\bigl(\varpi^{\sum_{s\in B_0}t_s}x\bigr)
 \in N(M^{\rm na})+\mathcal OC_\beta ,
\]
which is the all-but statement.

\paragraph{The level-index bound at the level prime.}
\teamanchor{proof.level-index-bound}%
Fix \(w\) and a place \(v\mid T_w\) of \(K\). Then
\(T_w\mathcal O_v=\mathfrak p_v^{e_v}\). Let \(\Lambda\) be the localization at
\(v\) of the zero-stage integral lattice of the column chain; this is
the lattice \(S_{0,w}\) above. It has
rank \(d_M\). Put \(\Lambda_r=\Lambda\cap V_r\). The release attaches to \(v\) a
finite set \(V^T_v\) of top vertices over \(T_w\) seen at \(v\). It has at
most \(N_T\) elements, and it is empty when \(v\) is not seen
\releasecite{Formalism/SharpRestriction/AdmissibleParameterJoins/DIVfromGLo.lean}{185}{The number of top vertices}. For each
\(s\in V^T_v\) the release supplies an endomorphism \(U_s\) of the rational
column space, the \(U_T\)-type Hecke operator at \(s\), with adjoint
\(U_s^*\). It also supplies a local lattice \(A_s\) and integers
\(a_s\), \(b_s\) and \(W_s\ge0\). The release inputs are as follows.
\begin{enumerate}
\item[(i)] The \(U_s\) commute. The operators \(U_s\) and \(U_s^*\) preserve
 \(V_r\), and \(U_sU_s^*=U_s^*U_s=T_w^{k-1}\) on \(V_r\). Moreover
 \(U_sA_s\subseteq\Lambda\) and \(U_s^*\Lambda\subseteq A_s\)
 \releasecite{Formalism/SharpRestriction/GlobalSchurBounds/ProductFirstZerothParts.lean}{195}{The Hecke relations over the level prime}
 \releasecite{Formalism/SharpRestriction/GlobalSchurBounds/FiniteModelOfAdmissibleTransports.lean}{2268}{Their proof from the admissible transports}.
\item[(ii)] \(b_s\ge0\), \(T_w^{b_s}\Lambda\subseteq A_s\), and
 \(\operatorname{length}(A_s/T_w^{b_s}\Lambda)=e_vd_M(a_s+kb_s)/k\)
 \releasecite{Formalism/SharpRestriction/GlobalSchurBounds/FiniteModelOfAdmissibleTransports.lean}{1557}{The auxiliary lattices at the vertices over the level prime}
 \releasecite{Formalism/SharpRestriction/GlobalSchurBounds/ProductFirstZerothParts.lean}{135}{The level-prime setting}.
\item[(iii)] \((a_s+kb_s)/k\le a_{\max}/k+2W_s\), where \(a_{\max}\ge0\) is the
 largest leading entry of the weights at the top places
 \releasecite{Formalism/SharpRestriction/GlobalSchurBounds/FiniteModelOfAdmissibleTransports.lean}{1706}{The twist inequality}.
\item[(iv)] \((N_w)_{(v)}\) is the preimage of
 \(\Lambda'_r=T_w^{e_T}U^m\Lambda_r\), with \(U=\prod_{s\in V^T_v}U_s\) and
 \(e_T=m\sum_sb_s\). The length of \((H_w)_v/(N_w)_v\) equals the length
 of \(\Lambda_r/\Lambda'_r\)
 \releasecite{Formalism/LocalRingCardinality/DataLatDefs.lean}{111}{The condition at the level prime}
 \releasecite{Formalism/SharpRestriction/AdmissibleParameterJoins/LenAtT.lean}{8}{The local length at the level prime}.
\item[(v)] \(a_{\max}\le2a_E\), \(W_s\le3b'\) and \(k\le b'\)
 \releasecite{Formalism/SharpRestriction/AdmissibleParameterJoins/DIVfromGLo.lean}{166}{The bound for the top weight sums}
 \releasecite{Formalism/SharpRestriction/AdmissibleParameterJoins/DIVfromGLo.lean}{112}{The bound for the weight widths}
 \releasecite{Formalism/SharpRestriction/AdmissibleParameterJoins/DIVfromGLo.lean}{44}{The rank parameters in the size data}.
\end{enumerate}
Put \(g_s=T_w^{b_s}U_s\). Then \(T_w^{b_s}\Lambda_r\subseteq A_s\cap V_r\) and
\(U_s(A_s\cap V_r)\subseteq \Lambda_r\), so \(g_s\Lambda_r\subseteq \Lambda_r\). Also
\(U_s^*\Lambda_r\subseteq A_s\cap V_r\), so
\(T_w^{k-1}\Lambda_r=U_sU_s^*\Lambda_r\subseteq U_s(A_s\cap V_r)\). The operator \(U_s\)
is injective on \(V_r\), because \(U_s^*U_s=T_w^{k-1}\) there. Hence
\[
 \begin{aligned}
 \operatorname{length}(\Lambda_r/g_s\Lambda_r)
 &=\operatorname{length}\bigl(\Lambda_r/U_s(A_s\cap V_r)\bigr)
  +\operatorname{length}\bigl((A_s\cap V_r)/T_w^{b_s}\Lambda_r\bigr)\\
 &\le\operatorname{length}(\Lambda_r/T_w^{k-1}\Lambda_r)
  +\operatorname{length}(A_s/T_w^{b_s}\Lambda)
 \le e_vd_M\Bigl(k-1+\frac{a_s+kb_s}{k}\Bigr).
 \end{aligned}
\]
The second term was bounded because \((A_s\cap V_r)/(T_w^{b_s}\Lambda\cap V_r)\)
embeds in \(A_s/T_w^{b_s}\Lambda\). If \(g\) and \(h\) are endomorphisms
preserving a lattice \(A\), then \(g\) maps \(A/hA\) onto \(gA/ghA\).
Therefore
\(\operatorname{length}(A/ghA)\le\operatorname{length}(A/gA)+\operatorname{length}(A/hA)\).
The \(g_s\) commute, and \(\Lambda'_r=(\prod_sg_s)^m\Lambda_r\). Using the
inputs (iii) and (v), we get
\[
 e_v\ell_{wv}=\operatorname{length}(\Lambda_r/\Lambda'_r)
 \le me_vd_M\sum_{s\in V^T_v}\Bigl(\frac{a_{\max}}{k}+2W_s+k\Bigr)
 \le7me_vd_MN_T\Bigl(\frac{a_E}{k}+b'\Bigr).
\]
In the level range \(2f/T_w\le\frac12\), we have
\(d_r\ge d_M/2\). Moreover
\(\sum_{v\mid T_w}e_v\lambda_v=\sum_{v\mid T_w}e_vf_v\log T_w/[K:\mathbb Q]=\log T_w\),
where \(f_v\) is the residue degree. Summing over \(v\mid T_w\) gives
\[
 \frac{2}{d_r}\sum_{v\mid T_w}e_v\lambda_v\ell_{wv}
 \le28N_Tm\Bigl(\frac{a_E}{k}+b'\Bigr)\log T_w .
\]
The release proves the bound in the same way, with its own constant
\releasecite{Formalism/SharpRestriction/AdmissibleParameterJoins/DIVfromGLd.lean}{70}{The level-index bound from its inputs}
\releasecite{Formalism/SharpRestriction/AdmissibleParameterJoins/DIVfromGLo.lean}{294}{Its inputs for the families}.

\paragraph{What the identification still takes from the release.}
The identification therefore reduces to these local
characterizations, which the release carries for every family in the
selected class. The edition derives the inclusion and the length
bound at served places from the release's definitions of the local
period data and from four release inputs. These are the
identification of the local lattice with the span of the source, the
lattice sandwich with its stable operators, the local lattice system
(i)--(iv) at the served vertices, and the rational spanning of the
local period lattice. The all-but statement itself is derived from
the lattice system by the walk argument. The edition derives the
level-index bound at \(v\mid T_w\) from the Hecke relations, the
auxiliary lattices, the twist inequality, the characterization at the
level prime with its length identity, and the size data listed there.
The edition still takes these inputs from the release, together with
the local characterizations themselves and the gain valuation on the
source at every place. The slope and norm deductions of
Section~\ref{sec:integral} use them as stated. The
\hyperlink{candidate.gain}{Frobenius divisibility and loss bound}
applies in its stated domain.

\paragraph{Which character is used at a final index?}
\teamanchor{proof.selected-ray-transport}%
The order of choices matters for the simultaneous specialization in
Section~\ref{sec:integral}. Use the family produced by the selection
above. Its preliminary index has the host, words, reference levels
and placed class data before the moving characters are chosen.
There is a finite set of marked places, defined by the row indices,
local shifts and content conditions. It is the inverse image of a
conjugation-stable set in a fixed coefficient subfield \(K_0\).
The Galois signatures, toric distributions and projected base
theta columns at this index determine the tests for specialization.
For each conjugacy orbit of marked places of \(K_0\) take one of the
ray characters furnished by simultaneous specialization. Carry the
previously chosen rational conductor primes into the next avoidance
set. There are finitely many orbits, so the supports for distinct
orbits can be disjoint. Keep the support, order and reciprocity
reading of each of these particular characters. At the other places
use the trivial character.
\releasecite{Development/AbelianVarietyTheory/Definitions/GlobalTwisting/DisjointConductorPlacesExcludedSet.lean}{139}{The three marking conditions}
\releasecite{Development/AbelianVarietyTheory/Proofs/IwasawaDecomposition/ExistsUnlocalizedSaturatedState.lean}{64}{The preliminary arithmetic data}
\releasecite{Development/AbelianVarietyTheory/Proofs/IwasawaDecomposition/TwistDataAgreeAcrossCoefficientsSymm.lean}{133}{The choice of ray data}.

The second pass of the family construction takes this finite list
as input. It adjoins its character values and the algebraic
constants in the test models to the large coefficient field. It
keeps \(K_0\) with its given embedding. The new marks are the
inverse images of the old ones. The base finite-level equations,
class representatives and real frames are carried to the same
base data. The construction realizes the moving characters with
the same ray groups, toric distributions and finite Hecke
correspondences. The final index retains the choice of side,
support and order. The geometric input here is the existence of
this matched adelic realization for the chosen preliminary data,
at every final index above one family threshold and before any
Fourier translate is picked.
\releasecite{Development/AutomorphicArithmetic/ThetaCorrespondence/ThetaPairLifts/ThirdPlaceSeries/FiniteLevelAndSecond.lean}{83}{Realization of the chosen characters}
\releasecite{Development/AbelianVarietyTheory/Definitions/AutomorphicMembership/CoefAEqOfHeq.lean}{166}{The compatible geometric enlargement}
\releasecite{Development/AbelianVarietyTheory/Proofs/IwasawaDecomposition/ExistsMemberOfConstructionStepsSecondPass.lean}{49}{Passage to the final member}.

\paragraph{The normalized theta series and its integral lattices under enlargement.}
\teamanchor{proof.selected-series-transport}%
Fix a test with the same stage, side, ray support \(Q\), order
and character in the two realizations. Write \(K\subset K'\)
for the large coefficient fields. The smaller fields and their
signed embeddings and local rings \(O_i,O_{i,k}\) are identified.
Let \(\omega\) be the Weil action and let \(\mathcal F(v,P)\)
be the Fock vector of the finite section \(v\) and harmonic
polynomial \(P\) in the prescribed infinite vacuum. Its theta
series at a finite source representative \(h\) and target point
\(g\) has value
\[
 T_v(h,P;g)=\bigl(\omega(h,g)\mathcal F(v,P)\bigr)(0).
\]
The matched lattices and split halves identify the base local
sections, the determinant test functions at \(Q\) and the
translations placing their product in the Weil model. The
splitting characters and harmonic frames agree, so the displayed
value agrees. For the universal series let \(v_a\) be the section
at ray coordinate \(a\), and let \(d_Q(g)\) be the ray coordinate
of the target determinant. Multiplication by its group-ring basis
element translates the coefficient at \(a\) to the value
\(T_{v_{a-d_Q(g)}}(h,P;g)\). Thus the coefficients already agree
before evaluation at any character. The specialized determinant
character on \(g\) also agrees. For a class form \(F\) the full
theta value is the sum of \(T_{v_\psi}(h,F(h);g)\) divided by
the order of the stabilizer of \(h\), over all source classes.
At corresponding levels the representatives and stabilizers agree
term by term.
\releasecite{Development/AbelianVarietyTheory/Definitions/NonGaloisDescent/AdmissibleEmptyOfGood.lean}{193}{The determinant shift}
\releasecite{Development/AbelianVarietyTheory/Definitions/NonGaloisDescent/AdmissibleEmptyOfGood.lean}{279}{The class sum}.

Use the same scalar \(\nu\) for the real normalization at this
stage and side. Indeed, transport the extremal vector and a
compatible polynomial frame. Their scale relation is unchanged
on a \(K\)-basis and hence on its \(K'\)-span. Changing between
two admissible frames of the irreducible lowest type multiplies
the map by a complex scalar \(\zeta\), by Schur's lemma applied
to their real models.
\releasecite{Development/AbelianVarietyTheory/Proofs/InertiaEigenvalues/StarMulSelfOfPair.lean}{24}{The scalar between two frames}
If the lowest type is zero take \(\zeta=1\). Otherwise let
\(\phi\) be the composite of one coefficient identification of the
frames with the inverse of the other. It is linear over the large
coefficient field, and the coefficient map \(c\) satisfies
\(c\circ\phi=\zeta c\). The frame data include that \(c\) sends a
basis of the coefficient space to complex linearly independent
vectors. Comparing coordinates of \(\phi\) in this basis, its
diagonal entries all have complex image \(\zeta\) and the others
vanish. Thus \(\phi\) is multiplication by an element \(a\) of the
large coefficient field with complex image \(\zeta\). We show that
\(a\) is a root of unity.
\begin{enumerate}
\item Let \(K_{\rm cm}\) be the common CM subfield, or the totally
real subfield if there is no CM subfield. The two identifications
send a fixed nonzero basis vector of the rational coefficient space
to \(ay\) and \(y\), and both have coordinates in \(K_{\rm cm}\) in
the fixed basis. Some coordinate \(y_k\) of \(y\) is nonzero, and
\(a=(ay)_k/y_k\in K_{\rm cm}\).
\item Both identifications carry the rational Hermitian pairing to
the same pairing \(\langle\,,\rangle\). Thus \(\langle y,y\rangle\)
is the rational pairing of the fixed basis vector with itself. Its
complex image is the squared norm of the image of that vector under
the injective complex coefficient map of the rational space, so
\(\langle y,y\rangle\ne0\). Hence
\(a\bar a\langle y,y\rangle=\langle ay,ay\rangle=\langle y,y\rangle\),
so \(a\bar a=1\).
\item Both identifications carry the integral harmonic lattice at a
class representative onto the same nonzero finitely generated
module \(\Lambda\) over the ring of integers of the large
coefficient field. Thus \(a\Lambda\subseteq\Lambda\). The module
\(\Lambda\) lies in a vector space, so it is faithful over the
ring generated by \(a\). By \cite[Prop.~5.1]{s2b-atiyah-macdonald}
the element \(a\) is integral over that ring of integers, hence
over \(\mathbf Z\) by transitivity \cite[Cor.~5.4]{s2b-atiyah-macdonald}.
\item Let \(\sigma\) be a complex embedding of \(K_{\rm cm}\). The
totally real subfield \(K_{\rm cm}^+\) fixed by the bar has real
image. If \(K_{\rm cm}\ne K_{\rm cm}^+\), write
\(K_{\rm cm}=K_{\rm cm}^+(\beta)\) with \(\beta^2\in K_{\rm cm}^+\)
and \(\bar\beta=-\beta\). Then \(\sigma(\beta)\) is not real,
because \(K_{\rm cm}\) is totally imaginary, while
\(\sigma(\beta)^2\) is real. So \(\sigma(\beta)^2<0\) and complex
conjugation sends \(\sigma(\beta)\) to \(-\sigma(\beta)\). It
therefore induces on \(\sigma(K_{\rm cm})\) the transport of the
bar. Thus
\(\sigma(\bar a)=\overline{\sigma(a)}\) in all cases, and
\(|\sigma(a)|^2=\sigma(a\bar a)=1\).
\end{enumerate}
The conjugates of \(a\) are the values \(\sigma(a)\). An algebraic
integer all of whose conjugates have absolute value one is a root
of unity \cite[Lemma~1.6]{i2-washington}. Hence \(a\) and
\(\zeta\) are roots of unity. The scale condition is agreement
of positive powers. Taking a multiple of the order of \(a\)
preserves it in both directions. On a tensor \(z\otimes f\) the
normalized evaluation is
\(\nu^{-1}\iota(z)\operatorname{ev}_{v_0}(f)(g)\), where
\(\iota\) is the signed embedding and \(v_0\) the extremal vector.
All three factors agree. Pure tensors span, so the rational
spaces of values and their membership conditions agree.
\releasecite{Development/AbelianVarietyTheory/Proofs/InertiaEigenvalues/ExtremalVectorMemLowestKType.lean}{234}{Admissibility of the same scalar}
\releasecite{Development/AbelianVarietyTheory/Proofs/IwasawaDecomposition/ExportTieMove.lean}{31}{The normalized evaluation}.

For the integral comparison, let \(J\) be the finite source
level with character \(\lambda_J\), and \(M_h\) the Hermitian
lattice obtained by translating the local source lattices by
\(h\). At \(h\) the coefficient functions of an integral
harmonic polynomial lie in the \(O_{i,k}\)-span of evaluations
of the dual Weyl tensors of \(M_h\). In matched local bases the
standard integral weight modules, their duals and the determinant
twists have the same matrices. Substitution in the polynomials
commutes with the identification of the small fields, hence these
integral harmonic lattices coincide. A class form must take
values in them and satisfy
\[
 F(\gamma h u)=\lambda_J(u)\rho(\gamma)F(h)
 \qquad(\gamma\in G(\mathbf Q),\ u\in J),
\]
where \(\rho\) is the polynomial action and \(G\) the source
group over \(\mathbf Q\). The stabilizer equations and this law
identify the integral forms in both directions. To pass back to
the coefficient coordinates, let \(m_j(g)\) be the finite matrix
of a geometric argument \(g\) at stage \(j\). The harmonic value
of a coefficient vector \(f\) at \(m_j(g)\) is the image of
the realized value \(f(g)\) under the fixed injective coefficient
map. The matched adelic matrices agree entry by entry, and the
finite representation and small-field action identify \(f(g)\).
The inverse images of the integral class forms under this map
are therefore the same at the two tests.
\releasecite{Development/AbelianVarietyTheory/Definitions/AutomorphicMembership/SubsetTwistExceptionalPlaces.lean}{123}{Harmonic evaluation on the finite representation}
\releasecite{Development/AbelianVarietyTheory/Proofs/IwasawaDecomposition/KeepsMatHeqRefl.lean}{22}{The finite matrix at this argument}.
In particular all
admissible base translates have the same support, coefficients
and theta values. Applying the same rational Hecke projection
and taking the span of their columns commutes with any further
extension of the coefficient DVR, even in residue characteristic
dividing the order of the level group. The universal coefficient
identity gives the corresponding assertion for the specialized
columns. The index comparing the base and the genuinely twisted
lattice at the final level is estimated in Section~\ref{sec:integral}.
\releasecite{Development/AbelianVarietyTheory/Definitions/NonGaloisDescent/AdmissibleEmptyOfGood.lean}{240}{Integral harmonic values and the level law}
\releasecite{Development/AbelianVarietyTheory/Proofs/InertiaEigenvalues/LiftSetupOverExtendedDataIntegralFormsEq.lean}{172}{Integral class coordinates in the two realizations}.

Assume the three estimates of Section~\ref{sec:integral} on the
preliminary signatures, toric distributions and theta tests.
To transfer them, take a marked place \(w\)
of the final coefficient field above \(\ell\), and let \(\alpha\)
be the preliminary character at its restriction.

\teamanchor{proof.selected-ray-frobenius}%
\paragraph{Identification of the final character.}
Let \(Q\) be the ray support of \(\alpha\) and \(G_Q\) its finite
ray quotient, with the continuous surjective Artin map
\(\operatorname{art}_Q\) from the absolute Galois group of \(E\).
The character \(\alpha\) takes values in the group \(\mu\) of
\(\ell^k\)-th roots of unity of its cyclotomic ring \(O_k\), which
contains a primitive \(\ell^k\)-th root. Let \(\iota_\alpha\) be
the fixed isomorphism of \(\mu\) with the complex roots of order
\(\ell^k\), and \(\chi\) the retained reciprocity reading. For
every place \(v\notin Q\) and every Frobenius element \(\sigma_v\)
at \(v\) it satisfies
\(\chi(v)^{-1}=\iota_\alpha\bigl(\alpha(\operatorname{art}_Q\sigma_v)\bigr)\).
Let \(\alpha'\colon G_Q\to O_k^\times\) be the final character at
\(w\), let \(L_w\) be the fraction field of the final coefficient
DVR at \(w\), and \(j_w\colon O_k\to L_w\) the structure map.
The final choice reads \(\alpha'\) through the same \(\chi\): if
\(v\notin Q\) does not lie above \(\ell\) and \(y\) in the large
coefficient field has complex image \(\chi(v)^{-1}\), then the
image of \(y\) in \(L_w\) is
\(j_w\bigl(\alpha'(\operatorname{art}_Q\sigma_v)\bigr)\).
The edition takes these two readings from the release: the
transport of \((Q,k,\alpha,\chi)\) to the final index with the same
ray quotient and Artin map, and the compatibility of the final
choice with \(\chi\).
\releasecite{Development/AbelianVarietyTheory/Proofs/IwasawaDecomposition/RawPackEqPow.lean}{51}{Transport of the preliminary ray data}
\releasecite{Formalism/LocalToGlobal/UnitaryNorms/CharGal.lean}{129}{Frobenius compatibility of the final choice}.

The second pass also makes the large coefficient field contain
every complex root of order \(\ell^k\); this is part of the
realization of the chosen characters taken from the release above.
The rest is derived. The complex embedding of the large coefficient
field is injective. Sending \(u\in\mu\) to the image in \(L_w\) of the
unique \(y\) with complex image \(\iota_\alpha(u)\) therefore
defines an injective homomorphism \(r_w\colon\mu\to L_w^\times\).
The restriction of \(j_w\) to \(\mu\) is also injective. The group
\(\mu\) is cyclic of order \(\ell^k\), and the \(\ell^k\)-th roots
of unity in the field \(L_w\) form a cyclic group of order at most
\(\ell^k\). Both maps are therefore isomorphisms onto this group,
and there is a unique \(u_w\in(\mathbf Z/\ell^k\mathbf Z)^\times\)
with \(r_w(u)=j_w(u^{u_w})\) for all \(u\in\mu\).
Now let \(\gamma\in G_Q\). The kernel of \(\operatorname{art}_Q\)
is open, so it fixes a finite abelian extension of \(E\) with
Galois group \(G_Q\). By Chebotarev's density theorem
\cite[Ch.~VII, Thm.~(13.4)]{i2-neukirch} the primes of \(E\)
unramified there with Frobenius \(\gamma\) have positive density.
Choose one, \(v\), outside \(Q\) and not above \(\ell\), and a
Frobenius element \(\sigma_v\) with
\(\operatorname{art}_Q\sigma_v=\gamma\). The two readings give
\[
 j_w\bigl(\alpha'(\gamma)\bigr)=r_w\bigl(\alpha(\gamma)\bigr)
   =j_w\bigl(\alpha(\gamma)^{u_w}\bigr).
\]
Since \(j_w\) is injective, \(\alpha'(\gamma)=\alpha(\gamma)^{u_w}\).
Hence \(\alpha'=\alpha^{u_w}\) on all of \(G_Q\). A primitive power on either prescribed side
is now a power \(\alpha^{\pm a u_w}\) with \(\ell\nmid a u_w\).
It has the same conductor and the same trivial restrictions at
the places over \(\ell\). The specialization in
Section~\ref{sec:integral} works for every such power on one
set of tests. When a coefficient DVR is extended, an integral
determinant valuation and the valuation of \(\ell\) are both
multiplied by its ramification degree. This preserves each
normalized signed lattice index and each toric valuation.
Extension is finite free on the Galois lattices. In an integral
basis it commutes with cohomology, the inverse image of the
rational finite condition and the kernels of the localization
maps. The resulting normalized Selmer lengths are also the same.
Thus the three estimates apply to the same character at this
index, with their original tolerances. The structural selection
identifies their arguments; their bounds come from the
corresponding local and global calculations.
\releasecite{Development/AbelianVarietyTheory/Proofs/IwasawaDecomposition/GoodTwistOfSeedCharacterPowers.lean}{157}{The same primitive powers in all the tests}.

\teamanchor{proof.same-family-hypotheses}%
Fix \(\varepsilon\) below \(1/52\), the existence cutoff and all the
comparison cutoffs. Apply selection once with \(K_*\). Use the same
\(C_H\ge1\) for the trace accuracy \(x^{-C_H}\) and the height test
\(H_*\le x^{C_H}b\). The following application uses the comparisons
in Sections~\ref{sec:family}--\ref{sec:fourier} on this family, including
the \hyperlink{proof.section-run-family-identification}{choice of
rows and field at each translated index} and the
\hyperlink{proof.ramified-donor-identification}{integral
identification of the ramified columns}.

\begin{itemize}
\item At a fixed sufficiently large \(x\in\operatorname{Win}(\mathfrak F)\)
choose \(i_0\in I_x\) and use its arrays as the seed. For each \(u\in H\)
choose \(i(u)\in I_x\) by the translate property. Its arrays agree with
the seed also for \(u=0\), whose index may differ from \(i_0\).
At \(i_0\), Section~\ref{sec:family} gives the three auxiliary quadratic
twist estimates at \(x\) and the tail estimate
\eqref{eq:host-character-tail} for every primitive associate of a host
character. At each \(i(u)\), the bound \eqref{eq:family-selmer} applies
to the selected field at its actual twist cutoff
\(t\le\lfloor x^{400}\rfloor\).

\item Use the matched period maps at \(i(u)\). The full base image and
small-level projections enter the rank comparison. For the actual
tuple at a finite place, \eqref{eq:period-selected-twist-avoidance}
gives the required prime avoidance, and
\eqref{eq:period-selected-local-row-lattice} identifies the localized
row lattice with the span of the projected integral forms.
The theta--Bessel normalization unit enters the lattice comparison, and the
Bessel maps into the corresponding Petersson space enter the Gram
comparison. At a nonbase twisted place the local Hecke ranks sum
to the dimension of the full definite source \(M^L_{\mu,j}\) in
\eqref{eq:record-local-rank-sum}. At the base twist the rank of
the good target \(M_j\) is \(d_j=D_0\) by
\eqref{eq:record-base-rank-sum}. It is this rank that occurs in
\eqref{eq:integral-selmer} for \(1\le j\le m\).
The lower bound \(m\ge x^{1/2-\varepsilon}\) and evenness give
\(m\ge4\) for large \(x\), so the bottom comparison has \(m/2-1\ge1\).

\item The period maps on the divided sources and into the one row
dual at \(i(u)\) carry the metrics of Sections~\ref{sec:trace} and
\ref{sec:integral}. The gain and lattice estimates use the Selmer
defect bound, the uniform Gram and reference-row bounds, the height
and slope comparisons in those metrics, and the centered level
count. On these hypotheses the lattice implication gives
\(D(u)\le E_x\) by \eqref{eq:integral-excess} for each \(u\), with
uniform constants.
\end{itemize}

At the same \(i(u)\) each \(h_w\) bounds the period map on its
divided source. The lower level bound and \(S_a=\sum_t a_t>0\)
give \(h_w\le h^*(u)\) for every word. Replacing each \(h_w\) by
this common upper bound lowers the adjusted mean to
\(s\le s_{\rm raw}\). Since the target is the row dual with its
squared Petersson norm, \(\mu_{\min}(C)=-\mu_{\max}(R)\).
The height cap \eqref{eq:raw-height-cap} therefore gives
\[
 \max(0,s-\mu_{\min}(C))\le x^{C_H}b .
\]
The common shift cancels in the difference of the donor slopes.
This proves the height condition for the frame estimate with the
exponent \(C_H\) already used in \eqref{eq:trace-accuracy}.

\paragraph{Two local rank comparisons.}
\label{sec:window-local-ranks}
\teamanchor{proof.window-local-ranks}
At a base-tuple stage, let \(L_j\) be the small-level isotypic
space and \(V_j\subseteq L_j\) its selected irreducible-type
part. The inclusion follows because the central type idempotent
commutes with the level idempotent, and \(V_j\) is the image of
their product. Let \(Z_j^{\rm coord}\subset V_j\) be the coordinate
row subspace used for the lifted tests. The actual block projection
from \(L_j\) to \(Z_j^{\rm coord}\) first projects to \(V_j\), then
projects for its positive Hermitian form onto \(Z_j^{\rm coord}\).
Its restriction to the included \(Z_j^{\rm coord}\) is the identity,
so it is onto. Indeed the first step is the product of the two
commuting idempotents. A product of commuting idempotents is an
idempotent, here with image \(V_j\), so it fixes \(V_j\) pointwise.
The Hermitian form is positive definite on \(V_j\), so
\(V_j=Z_j^{\rm coord}\oplus(Z_j^{\rm coord})^\perp\) and the
orthogonal projection fixes \(Z_j^{\rm coord}\) pointwise. Thus
projection after inclusion is the identity of \(Z_j^{\rm coord}\).
This is the right-inverse identity for the two small-level maps in
the projected coordinates.

Put \(i=j-1\). The map entering this rank is the auxiliary lift
\(\Theta_{P_0,\boldsymbol\mu,j}\) from the definite group of rank \(r+t\) to
the group of rank \(r\). The descent through the spaces \(Y_k\)
of ranks \(r+kt\) supplies the bottom image used in the period
comparison. For the auxiliary lift use the source class-coordinate
inclusion \(I_{i,\boldsymbol\mu}\), the fitted projection \(P_ZV^\flat E_J\),
and the target basis map \(j_{W,i}\) of
\eqref{eq:period-fitted-stage-vector}. The one-step norm of
Section~\ref{sec:stage-petersson} is positive on each occurring
target summand. Applying the commuting good and level projectors
to a preimage gives surjectivity on the cut. Thus the restriction
to the rational kernel complement has invertible matrix
\(U_{i,\boldsymbol\mu}\). For a standard target coordinate \(e_a\), set
\[
 x_a=E_J^\flat VZ\,U_{i,\boldsymbol\mu}^{-1}e_a.
\]
Applying the right-inverse identity and then the column formula gives
\[
  (P_ZV^\flat E_J)x_a=U_{i,\boldsymbol\mu}^{-1}e_a,\qquad
  \Theta_{P_0,\boldsymbol\mu,j}(I_{i,\boldsymbol\mu}x_a)=j_{W,i}(e_a).
\]
The vectors on the right span the stage multiplicity space \(M_j\).
The kernel of the adjoint \(\Theta_{T,j}\) is orthogonal to their
span, so the adjoint is injective.

\paragraph{The base record and its finite-type blocks.}
\teamanchor{proof.base-record-blocks}%
Use the common small level \((P_0,\psi_0)\) and base tuple
\(\boldsymbol\mu_0\). The bottom row space is the selected space on
\(U(Y_0)\), with the real model used for the bottom period. For
\(j=i+1\), the space \(M_j\) in the record is read on the auxiliary
rank-\(r\) group \(U(W_i)\). Its fitted theta map comes from the
rank-\((r+t)\) auxiliary group \(U(V_i)\), as in
\eqref{eq:period-fitted-stage-vector}. The descent from the definite
top of the \(Y\)-tower enters through the bottom period and its
spherical eigenvalue transfer.

First compare the finite group actions. Let \(\Gamma\) be the
product of the rank-\(r\) residue groups at the type places. The
principal level is the kernel of the rank-\(r\) reduction within
the subgroup acting in each selected frame. The suffix character
is included in the lift of this action. In the adelic realization
one takes right translation by a lift of \(g\in\Gamma\), divided
by its prescribed suffix character. The covariance under the
principal level shows that the result is independent of the lift.
For a definite class value, write the translated representative as
\(\gamma g_c u\), with \(\gamma\) rational and \(u\) in the level
subgroup. The resulting value is the same coefficient action of
\(\gamma\) on the value at \(g_c\), with the level factor from
\(u\). This is the class-coordinate matrix \(\rho(g)\) of the
action. For a character \(\eta\) of \(P_0\), a central element
\(z=\sum_g z_g g\) acts on its level coordinates by
\[
 E_{P_0}^{\flat}\Bigl(\sum_g z_g\rho(g)\Bigr)E_{P_0},
 \qquad
 e_{P_0,\eta}=|P_0|^{-1}\sum_{h\in P_0}\eta(h)^{-1}\rho(h).
\]
Here \(E_{P_0}\) and \(E_{P_0}^{\flat}\) are the prescribed
inclusion and orthogonal coordinate projection, whose composite
in the other order is the level projector displayed on the right.
For both normalized actions in the fitted comparison the residue
character is \(\eta=\psi_0^{-1}\). On the auxiliary source this comes
from the ratio of the enlarged level character to the source flag
character. On the row it is the level character of the read model.
Thus the coefficient of \(\rho(h)\) in the level average is
\(\psi_0(h)/|P_0|\) on both sides.
Centrality gives commutation with that projector. Taking \(z\) to
be the sum of the central idempotents of the allowed Green types
gives the prescribed good cut. The same averaging acts on the
actual automorphic vectors. Complex conjugation between the two
bottom realizations inverts the finite level character and
conjugates the central projector.
The allowed list is stable under duality, so the central projector
is unchanged by conjugation. The scalar \(\sqrt{n_{P_0}}\) in the
level embedding of \(M_j\) cancels when transporting the group
actions.
\releasecite{Formalism/LocalToGlobal/Divisibility/RowLat.lean}{248}{The inverse character in the auxiliary projector}
\releasecite{Formalism/LocalToGlobal/Divisibility/GoodPartHermitianSpec.lean}{116}{Level coordinates and central matrices}
\releasecite{Formalism/LocalToGlobal/HilbertRowConditions/CenterIdempotentAction.lean}{17}{The same claim for the prescribed bottom level}
\releasecite{Development/AbelianVarietyTheory/Definitions/AutomorphicMembership/StableIffCarrierEq.lean}{116}{The level embeddings and their normalizations}.

Now let \(\sigma\) be a finite adelic isomorphism class occurring
in either rank-\(r\) good space. It contains a tensor of allowed
finite types in its principal-level fixed space. Indeed, the
good projector is a central sum of isotypic projectors, and a
nonzero vector of its image has a nonzero component for one
summand. At the base level \(P_0\), the factors of \(\psi_0\) are
nondegenerate Whittaker characters of the chosen maximal unipotent
groups. Every allowed tensor contains such a line. In each local
parabolic induction the open parabolic--unipotent cell gives the
tensor product of the cuspidal Whittaker lines
\cite{tr3green,tr3gelfand}. Conjugation by the alternating diagonal
inverts the character. A type occurring at either of the two bottom
levels therefore occurs at the base level in the same finite adelic
representation. The lifts of the residue action preserve the common
principal level and all level conditions away from the type places.
\releasecite{Formalism/LocalToGlobal/IndefiniteReal/LinksDefs.lean}{430}{Level away from the type groups}
\releasecite{Development/AbelianVarietyTheory/Proofs/SatakeParameters/GoodGenericCharacterOfSeriesTypes.lean}{23}{The two signs on the Whittaker line} and
\releasecite{Development/AbelianVarietyTheory/Proofs/TwistedOperators/BlockFrameConjugationIdentity.lean}{243}{The resulting base vector}.

By the Green parametrization \cite{tr3green} each allowed local factor
is induced from the distinct cuspidal factors associated with
the irreducible factors of its squarefree polynomial. If an
\((a,r-a)\) unipotent radical fixed a nonzero vector, the finite
Mackey formula would admit a double coset for which each cuspidal
block lies entirely on one side of the cut. A block meeting
both sides would give a proper unipotent radical of that block
and its cuspidality gives zero. Thus \(a\) must be a subsum of
the factor degrees. The first type-prime condition excludes a
common subsum of degree \(a\) at all the selected places over
\(g_0\) for \(1\le a\le r/2\). Taking complements gives the
exclusion for \(r/2<a<r\). For each \(a\), therefore, some actual
type factor in \(\sigma\) has zero invariants under that radical.
The criterion also holds for the dual types.
\releasecite{Development/AbelianVarietyTheory/Proofs/TwistedOperators/IrreducibleNotKilledByGoodProjectorSatisfiesH.lean}{159}{The Green test for a surviving type}
\releasecite{Development/AbelianVarietyTheory/Proofs/DivisorProofs/ExistsTypeVectorOfGoodLevelRepresentation.lean}{116}{Insertion into the occurring finite representation}.

The same test excludes a global proper block of size \(a\).
At the chosen split place, the group is \(\mathrm{GL}_r\) and
the principal compact subgroup is \(K(1)\). By Iwasawa
decomposition, the \(K(1)\)-invariants of normalized induction
from the \((a,r-a)\) parabolic are the finite parabolic
induction of the inducing Levi's principal invariants. Taking
these invariants is exact by compact averaging, and finite
Frobenius reciprocity shows that each irreducible finite type
in a subquotient has nonzero invariants under the finite
unipotent radical. The same conclusion applies to a refinement
of the parabolic by induction in stages. This contradicts the
type factor just found. In particular, the discrete automorphic
occurrences with this finite cut are cuspidal, since a residual
constituent is a subquotient of induction from a proper global
Levi and its block sizes specialize at each split place.

We compare multiplicities for this particular \(\sigma\).
The two rank-\(r\) groups are Hermitian pure inner forms of the
same quasi-split unitary group. Their matched finite lattice
frames identify the local groups by inner isomorphisms, carrying
the level subgroup, its character and the residue homomorphisms.
At all real places the conjugate bottom model and the model on
\(W_i\) have the same list of Harish--Chandra parameters with
signs forgotten. The selected signatures and assignments of
the entries to the two signs determine a discrete-series member
on each group. The lists are integral at this odd rank, with
adjacent gaps greater than \(1\).
\releasecite{Development/AbelianVarietyTheory/Proofs/DivisorProofs/TraceEqXOfGeomExpansionX.lean}{502}{The real weights and their gaps}
\releasecite{Development/AbelianVarietyTheory/Proofs/IwasawaDecomposition/FinrankBaseRowsIsotypicEqGoodFormsIsotypic.lean}{97}{Matched frames on the same finite isotypic space}.

For either actual occurrence, Labesse's weak base change
identifies local base change at every split place as well as
at the unramified places. The hypotheses on the CM extension
and real field hold for the chosen host, including
\([F:\mathbf Q]\ge 2\). The archimedean sign condition in that
theorem follows from the gaps. In fact Salamanca--Riba's
classification, applied to \(\mathrm{SU}(p,q)\) with the central
character fixed on \(U(p,q)\), makes a unitary representation
of this infinitesimal character an \(A_{\mathfrak q}(\lambda)\)
induced from a character of a \(\theta\)-stable Levi. In any Levi
block of size greater than one, both \(\lambda\) and
\(\rho(\mathfrak u)\) are constant on the block and
\(\rho(\mathfrak l)\) has two entries at distance one. The
infinitesimal character would contain two entries at distance
one. Hence the Levi is toral, giving a discrete-series member
of the same real packet. These are precisely the representations
allowed by Labesse's criterion for the pseudo-coefficient at
infinity \cite{rank-labesse,rank-salamanca}.

Write the base change as an isobaric sum of discrete
representations of general linear groups. With several summands,
the rank \(a\) of one summand yields the proper parabolic
test above at each split place. With a single noncuspidal
discrete summand, the M\oe glin--Waldspurger description is
Speh from a fixed cuspidal rank \(a<r\). Its local components
are subquotients of induction of twists of that same rank,
so the same test applies. Therefore the stable base change is
a single cuspidal \(\Pi\) of \(\mathrm{GL}_r(\mathbb A_E)\).
The comparisons here use Labesse's local statement at the
ramified split type places as well \cite{rank-labesse,rank-moeglin}.

In the Kaletha--M\'inguez--Shin--White classification take
\(\kappa=+1\), \(\chi_+=1\), which gives this same stable
base change. Their Corollary~3.3.2 attaches the Satake records
of the actual discrete occurrence to a global parameter.
At almost every place of \(E\) its general-linear record
equals that of \(\Pi\), by local base change at split and
inert places. Theorem~4.4 of Jacquet and Shalika identifies the
cuspidal inducing factors, up to permutation, from this equality of
unramified components outside a finite set \cite{rank-jacquet-shalika}.
There is just one factor, namely \(\Pi\), of rank \(r\) and with
Arthur \(\mathrm{SL}_2\) of dimension one. For a simple
parameter the centralizer modulo the invariant dual center is
trivial. Apply the proved generic, pure-inner case of
Theorem*~1.7.1. The transported finite factors stay in the
same local packets by the inner isomorphisms. At infinity
the common unsigned list gives the same local discrete
parameter and the chosen signs give the specified member on
the other group. With trivial global component group this
tensor therefore occurs there with multiplicity one. Starting
on either group gives the same argument \cite{rank-kmsw}.

Consequently the contribution of \(\sigma\) to either level is
the dimension of the same finite space
\(e_{\mathrm{good}}\sigma^{J_0,\psi_0^{-1}}\), where \(J_0\) denotes
the transported level. Use its conjugate realization on the
bottom. The record's \(D_0\) is the size of the chosen basis
\((b_k)\) of that bottom space with type labels \(\operatorname{ft}_0(k)\).
The stage labels are read by the actual finite isotypic spaces
on \(W_i\). In particular, if \(M_j[\sigma]\) denotes that
stage isotype, then at each labelled \(\sigma\) the identity is
\begin{equation}
 \label{eq:record-individual-isot-rank}
 \dim_{\mathbf C} M_j[\sigma]
 = \#\{k:\operatorname{ft}_0(k)=\sigma\}.
\end{equation}
The matched groups also give the reverse inclusion on occurring
labels, so these spaces exhaust \(M_j\).
\releasecite{Development/AbelianVarietyTheory/Definitions/AutomorphicMembership/RowsRestWordColumnUnitBaseTupleAlt.lean}{356}{The actual basis, finite-type labels and auxiliary stage}
\releasecite{Development/AbelianVarietyTheory/Proofs/IwasawaDecomposition/FinrankIsotEqIsotL.lean}{17}{Reading the stage isotype in its original frame}.

Finally match the spherical labels. Off the type and exceptional
places, each finite representation \(\sigma\) has a spherical
line. Its scalar for the double coset of weight \(d\) on
\(W_i\), computed in the matched lattice frame, is denoted
\(c_\sigma(v,d)\). Inner finite transport makes this independent
of \(i\). For the conjugate bottom representation the scalar is
\(\overline{c_\sigma(v,d)}=c_\sigma(v,d^\vee)\), with
\(d^\vee=(-d_r,\ldots,-d_1)\). The equality follows by taking
the adjoint, which is the inverse double coset. This accounts for
the dual convention in the bottom eigenspace. At inert places
the same argument uses the relative Cartan double cosets.

The top base record \(p\) consists of the prescribed split
Hecke polynomials and the inert double-coset scalars of the
definite top space. The bottom space for \(c_\sigma\) is carried
by \(p\) when its long-chain adjoint lies in the top block
\(p\). On each chosen split and inert spherical operator, the
adjoint carries the eigencharacter by the local spherical theta
transfer. At each arrow the unitary Satake list has the old
entries with the chosen common character twist and the prescribed
new string, with the dual convention just described if using the
conjugate model. The spherical calculation applies to the
hyperspecial lattice vectors and unramified characters at these
places. For an inert place orient the correspondence from the
smaller group to the larger one. In Harris's notation this is
\(m\ge n\). The lattice vector gives the spherical line, and the
map of spherical parameters on the domain is a morphism. Its
agreement with the string formula on a dense open set consequently
holds for each parameter that occurs here. The additive and
splitting characters are unramified and the lattices self-dual.
The residue characteristic is odd since the places over \(2\)
split in the chosen host. The finite level
and central eigencuts commute with these operators, which act at
disjoint places. The long-chain adjoint is injective on the
selected bottom space by the base row descent, so its nonzero
image belongs to the unique top block with these joint scalars.
The rational top blocks are the joint kernels of the same
operators with those eigenvalues. Their base changes are the
corresponding complex kernels by the class-coordinate comparison.
Thus the assigned label \(p(\sigma)\) is an actual index of the
prescribed top record. Conversely this index gives the full split
Hecke polynomials and the scalars of the inert Cartan generators.
At an inert place the relative Satake class is determined by its
reciprocal pairs, with middle entry \(1\) in odd rank. The relative
Weyl group permutes the pairs and interchanges their two entries.
At each lift remove the prescribed string with its multiplicities
and undo the common character twist. This recovers the spherical
character of the bottom representation on every operator in the
record. Several finite representations may have the same label
\(p(\sigma)\), and the spherical character on that block is
well defined \cite{rank-harris-spherical}.
\releasecite{Formalism/LocalToGlobal/IndefiniteReal/ZeroPartEigenvalueConvention.lean}{103}{The dual bottom convention and the carried label}
\releasecite{Formalism/LocalToGlobal/Divisibility/InertRationalAdjoint.lean}{22}{The prescribed split and inert index}
\releasecite{Development/AbelianVarietyTheory/Proofs/IwasawaDecomposition/HasLab.lean}{16}{Reading the label on the actual vectors} and
\releasecite{Development/AbelianVarietyTheory/Proofs/TwistedOperators/FiniteTypeCarriesOnePart.lean}{91}{Recovering the spherical summand}.

The bottom coordinate summands \(E_p\) are the preimages of
the corresponding analytic cuts under the base theta columns.
The top block projectors are defined over \(K\). The long-chain
adjoint maps each bottom eigenspace into a single top block, so
the projectors preserve the image of the adjoint, as well as the
rational kernel complement whose scalar extension is this image.
Taking adjoints shows that each rational summand maps into its
indicated analytic cut. These summands span the complement,
the map is injective there, and its complexification has image
the whole selected bottom space by the base row descent.
Orthogonality of distinct spherical parts gives surjectivity
block by block. Similarly the auxiliary coefficient summand
\(W_{i,p}\) is the preimage of its analytic cut under the
class-basis embedding of \(W_i\). The base-tuple complementary
stage coordinates split by the same transported top labels.
Their images under the fitted matrix belong to the corresponding
analytic cuts by \eqref{eq:period-fitted-stage-vector} and the
auxiliary spherical transfer. These coordinate summands exhaust
the complementary coordinates, and the base matrix on the
complement is invertible. Thus their images span the full good
class-basis space over \(K\) and split it by the given labels.
The prescribed good basis and its scaled level reading give the
complexification of each such summand. Equivalently, these are the
direct sums of the actual isotypes with \(p(\sigma)=p\). The
finite isotypic summands are orthogonal because they are distinct
irreducible representations of the finite adelic group in the
unitary discrete spectrum.
\releasecite{Formalism/LocalToGlobal/IndefiniteReal/PartsOrth.lean}{137}{Definitions of the rational preimages}
\releasecite{Formalism/LocalToGlobal/IndefiniteReal/FinrankISupLeSum.lean}{116}{Passage to the same individual parts}
\releasecite{Development/AbelianVarietyTheory/Proofs/IwasawaDecomposition/ComplexSpanRationalLiftFormsEqBasePart.lean}{22}{The good basis on the auxiliary part}.

It follows from \eqref{eq:record-individual-isot-rank} that
\begin{equation}
 \label{eq:record-base-rank-sum}
 \begin{aligned}
 \dim_K E_p=\dim_K W_{i,p}
   &=\sum_{\sigma:p(\sigma)=p}\#\{k:\operatorname{ft}_0(k)=\sigma\},\\
 \sum_p\dim_K E_p&=D_0=\dim_{\mathbf C} M_j .
 \end{aligned}
\end{equation}
\teamanchor{proof.base-record-rank-count}%
In detail, dimension is unchanged by the scalar extension from
\(K\) to \(\mathbf C\). By the preceding paragraph the
complexification of \(E_p\) is the direct sum of the bottom
isotypes of the \(\sigma\) with \(p(\sigma)=p\), and that of
\(W_{i,p}\) is the direct sum of the stage isotypes
\(M_j[\sigma]\) with \(p(\sigma)=p\). The chosen labelled basis
\((b_k)\) is adapted to the bottom isotypic decomposition: for
each \(\sigma\) the \(b_k\) with \(\operatorname{ft}_0(k)=\sigma\)
span the bottom \(\sigma\)-isotype. This is part of the release's
description of the labelled basis cited above. So the bottom
\(\sigma\)-isotype has dimension
\(\#\{k:\operatorname{ft}_0(k)=\sigma\}\), and by
\eqref{eq:record-individual-isot-rank} so has \(M_j[\sigma]\).
Summing over the \(\sigma\) with \(p(\sigma)=p\) gives the first
line. Each occurring \(\sigma\) has exactly one label \(p(\sigma)\),
and each \(k\) has the occurring label \(\operatorname{ft}_0(k)\).
Hence the double sum over \(p\) and \(\sigma\) counts every
\(k\) once, and \(\sum_p\dim_KE_p=D_0\). Since
\(M_j=\bigoplus_\sigma M_j[\sigma]\), also
\(\dim_{\mathbf C}M_j=\sum_\sigma\#\{k:\operatorname{ft}_0(k)=\sigma\}
=D_0\).
This is the rank of the indicated base row space on the same
level. With the standard-vector equality and the level coordinate
right inverse, the fitted map onto this \(M_j\) is surjective,
by the vectors \(x_a\) above. For its Hilbert adjoint,
\(\ker\Theta^*=(\operatorname{im}\Theta)^\perp=0\) and
\(\operatorname{rank}\Theta^*=\operatorname{rank}\Theta
=\dim M_j=D_0\).

On a nonzero carried block, the auxiliary Petersson scalar in
\eqref{eq:actual-stage-petersson-scalar} is the scalar used when
constructing the record at the distinguished embedding. Indeed,
the scalar is chosen there by the equation
\(\Theta\Theta^* y=\lambda_{j,\mu,p}y\) for the same phase-one
auxiliary theta map and the same carried space. On this block
its value is determined by the ratio
\(\|\Theta^*y\|^2/\|y\|^2\) for any nonzero \(y\). The
level-index scaling of the two Petersson forms is the same
and therefore cancels in the ratio. The normalized comparison
\eqref{eq:period-fitted-stage-vector} gives the same ratio for
the fixed record map and its adjoint. Consequently, for each
used tuple and each nonzero \(M_j[\sigma]\), the record scalar
\(\lambda_{j,\mu,\sigma}\) equals \(\lambda_{j,\mu,p(\sigma)}\)
at this embedding.
\releasecite{Formalism/LocalToGlobal/IndefiniteReal/WeakCharacterisationWithError.lean}{111}{Selection of the actual eigenvalue}
\releasecite{Formalism/LocalToGlobal/IndefiniteReal/BaseLiftsPackageOfPrimitiveLinks.lean}{68}{Auxiliary theta and its Petersson adjoint}
\releasecite{Development/AbelianVarietyTheory/Proofs/IwasawaDecomposition/LiftScalarEqLabelScalar.lean}{28}{The scalar on a nonzero finite isotype}.

\paragraph{The residual Hecke summands at a nonbase twist.}
\teamanchor{proof.nonbase-record-blocks}%
Fix the stage \(j\) and coefficient place \(w\) of a nonbase twist
\(\mu=\mu_{\rm tw}(w)\). Let \(M^L_{\mu,j}\) be the full rational
space of definite forms of rank \(r+t\) at the prescribed level
character. It uses the class representatives and coefficient
representation of \eqref{eq:period-auxiliary-record-read}.
The good-type cut will be taken inside the resulting local lattices.
Write \(O\) for the local coefficient DVR, \(L\) for its fraction
field, \(M=L\otimes_K M^L_{\mu,j}\),
and \(A\subseteq\operatorname{End}_L(M)\) for the \(O\)-algebra
generated by the prescribed normalized Hecke operators. To specify
them, put \(n=r+t\) and \(\ell=\operatorname{char}k(w)\). Each
generator is indexed by a split place \(v\) of \(E\) away from
\(\ell\), the two type primes, the designated scalar-character
prime and the conductor set of the twist. The last set includes
both halves above each place of \(F\) in its support.
Use the split realization at \(\bar v\), put \(q=|\mathcal O_E/v|\),
and write \(J\) for the hyperspecial group in those coordinates.
For \(1\le k\le n\) its action on the definite class forms is
\begin{equation}
 \label{eq:record-normalized-hecke}
 h_{v,k}f(g)=q^{-k(n-1)}
    \sum_{a\in J\operatorname{diag}(\varpi I_k,I_{n-k})J/J}
                   f(ga).
\end{equation}
Here \(\varpi\) is a uniformizer at the underlying place of \(F\).
Using the conjugate half fixes which eigenvalues occur in this
formula. In the normalized Satake convention with parameters
\(\alpha_1,\ldots,\alpha_n\) at that half, the coset sum has
eigenvalue \(q^{k(n-k)/2}e_k(\alpha_1,\ldots,\alpha_n)\).
Indeed the right cosets correspond to sublattices between
\(\varpi\mathcal O^n\) and \(\mathcal O^n\) of residual codimension
\(k\). In upper triangular echelon form fix the \(k\) positions
\(I\) with diagonal \(\varpi\). There is one free residue for each
\(a\in I\), \(b\notin I\) with \(a<b\), hence
\(\sum_{a\in I}(n-a)-k(k-1)/2\) free residues. On this cell
the normalized spherical vector takes the value
\(q^{-\sum_{a\in I}(n+1-2a)/2}\prod_{a\in I}\alpha_a\).
The cell size times the power of \(q\) is \(q^{k(n-k)/2}\),
independently of \(I\). Summing proves the asserted eigenvalue.
Consequently the monic polynomial with coefficient
\((-1)^kq^{k(k-1)/2}h_{v,k}\) in degree \(n-k\) has roots
\(q^{(1-n)/2}\alpha_a\) on this eigenspace. At the indexed half
\(v\) the Satake list is the inverse list, by the inverse-transpose
identification of the split unitary group. Arithmetic Frobenius in
the convention of the record inverts that list once more before
applying the norm shift.
\releasecite{Development/AbelianVarietyTheory/Proofs/IwasawaDecomposition/SphericalVectorFixedLocal.lean}{96}{The half-place and Frobenius conventions}.

Each representative and its inverse have denominators supported at
the rational prime below \(v\). Thus the two value lattices at \(g\)
and \(ga\) have the same localization at \(w\), including for the
coefficient representation. The finite sum preserves the localized
class lattice, and \(q\) is a unit there. Its value at an adelic
argument is reconstructed from the rational coefficient action and the
level character on its class representative. This puts the coset
operator in the prescribed rational class model. The resulting operators
have integral characteristic polynomials. They commute by the
spherical Hecke algebra at each place and independence across
places. The \(L\) used here already contains their eigenvalues:
the commuting endomorphisms of the finite-dimensional class space
generate a finite-dimensional algebra, so only finitely many of
them are needed to generate it over the fraction field. One may
split those generators in a finite extension. The local coefficient
choice at \((j,w)\) takes the compositum of that extension with the
prescribed coefficient and twist fields before forming the residual
ideals. The split commuting generators can be triangularized
simultaneously, so their algebra splits over this one \(L\).
Its generator eigenvalues belong to the valuation ring \(O\).
\releasecite{Development/AbelianVarietyTheory/Proofs/IwasawaDecomposition/ExistsConjugateSplitPlaceAndUniformizer.lean}{10}{Choice of the split half}
\releasecite{Development/AbelianVarietyTheory/Proofs/IwasawaDecomposition/HeckeOperatorKeepsRationalIntegralLevelForms.lean}{203}{Localized coefficient lattices under the coset sum}
\releasecite{Development/AbelianVarietyTheory/Proofs/IwasawaDecomposition/SideRationalEigenData.lean}{124}{The common splitting field at the coefficient place}.

Successive primary decompositions give joint generalized
eigenspaces. Indeed each generator preserves the spaces already
obtained, and each proper refinement increases the number of
nonzero summands, which is at most \(\dim_L M\). On a joint summand
the commuting nilpotent parts generate a nilpotent ideal in the
finite-dimensional commutative image algebra. Taking the scalar
part therefore gives a character \(A\to O\). Its reduction modulo
the uniformizer is onto the residue field since \(A\) contains \(O\).

\teamanchor{proof.record-maximal-ideals}%
Let \(\theta_1,\ldots,\theta_s\) be the distinct characters so
obtained and \(M_i\) the joint summand of \(\theta_i\). Then
\(\mathfrak m_i=\theta_i^{-1}(\varpi_wO)\), with \(\varpi_w\) a
uniformizer of \(O\), is a maximal ideal of \(A\), being the kernel
of a surjection onto the residue field. Every maximal ideal of
\(A\) is one of the \(\mathfrak m_i\). Indeed, let
\(\theta=(\theta_1,\ldots,\theta_s)\colon A\to O^s\). An element of
its kernel acts nilpotently on each \(M_i\), hence on \(M\). Since
\(A\subseteq\operatorname{End}_L(M)\), this kernel is a nil ideal
and lies in every maximal ideal. Each generator has an integral
characteristic polynomial, so it is integral over \(O\) by the
Cayley--Hamilton theorem. The elements of the commutative ring \(A\)
integral over \(O\) form a subring \cite[Cor.~5.3]{s2b-atiyah-macdonald},
which contains the generators and \(O\). Hence \(A\) and its image
\(B=\theta(A)\) are integral over \(O\). Also \(O^s\) is integral
over \(B\): an element \(x\) is a root of \(\prod_i(X-x_i)\), whose
coefficients lie in the diagonal copy of \(O\) in \(B\). Let \(\mathfrak m\) be a maximal ideal of \(A\). Its
image is a maximal ideal of \(B\). By lying over and
\cite[Thm.~5.10 and Cor.~5.8]{s2b-atiyah-macdonald} it is the
contraction of a maximal ideal of \(O^s\). The maximal ideals of
\(O^s\) are the inverse images of \(\varpi_wO\) under the
coordinate projections. Hence \(\mathfrak m=\mathfrak m_i\) for
some \(i\). In particular \(A\) has finitely many maximal ideals.

For a maximal ideal \(\mathfrak m\) of \(A\), let \(M_{\mathfrak m}\)
be the sum of the \(M_i\) with \(\mathfrak m_i=\mathfrak m\) and
\(p_{\mathfrak m}\) the projection along the other summands. The
local indices over \((j,w)\) run once through the maximal ideals of
this algebra; the record carries this indexing and the lattices
below.
\releasecite{Formalism/LocalToGlobal/NumberFieldFamilies/NPartCNd.lean}{132}{The local indices and the projected generator lattice} For such an index the lattice is the \(O\)-span of
\(p_{\mathfrak m}b\), where \(b\) runs through the realized generators
of the specified twisted class lattice. These generators span \(M\)
over \(L\). Applying \(p_{\mathfrak m}\) shows that the local
lattice spans exactly \(M_{\mathfrak m}\) over \(L\). It is finite
free, being a finitely generated torsion-free module over a DVR.
Its rank is \(\dim_LM_{\mathfrak m}\). Each \(M_i\) occurs in
exactly one \(M_{\mathfrak m}\), and \(M\) is the direct sum of the
\(M_i\). Consequently, writing \(\Lambda_a\) for the lattice with local
index \(a\), we have
\begin{equation}
 \label{eq:record-local-rank-sum}
 \sum_{a\in\operatorname{Loc}(j,w)}\operatorname{rank}_O\Lambda_a
   =\sum_{\mathfrak m}\dim_L M_{\mathfrak m}
   =\dim_L M=\dim_K M^L_{\mu,j}.
\end{equation}
If \(M=0\), its lattice is zero and the index set is empty.
In every case the local projections jointly detect each vector of
\(M\), since their sum is the identity. This rank computation takes place
over the fraction field. The colength in
Section~\ref{sec:integral} also depends on the position of the
projected integral lattice relative to the intersected lattice.

The two estimates use the same \(D(u)\). Indeed \(T,T_2,a_t\) agree
by the translate property, the weights \(w_{u,j}\) in both sums
come from \(i(u)\), and the inertia classes and sign character are
read at \(i_0\), with the ramified inertia groups of
\eqref{eq:classical-host-inertia}. By
\eqref{eq:prime-power-cutoff} every positive-part summand with
\(p^h>N=\lceil2xb\rceil\) vanishes. Thus the cutoff and the
untruncated Fourier ramp give the same \(A(a_t,t,w_{u,j})\).
Both comparisons subtract the real part of \(M(a_t,w_{u,j})\) in
\eqref{eq:prime-power-center}, with \(\gamma\) and the sign-character
log derivative at \(1\) in the same normalization. Hence the Fourier
estimate uses exactly \eqref{eq:centered-excess-definition}.

With the family, precision and constants fixed, take \(x_*\) to be
the maximum of the family, arithmetic, period, trace, lattice and
Fourier thresholds, increasing it to \(3\) if necessary. Each
threshold in this maximum works for every required index at the
scale. At every \(x\ge x_*\) in \(\operatorname{Win}(\mathfrak F)\)
the preceding applications give the lattice bound for all \(u\) and the
base-host character estimates. The sampling and weight comparisons
of Section~\ref{sec:fourier} refer to these same records. Thus
\eqref{eq:fourier-pointwise} applies with \(e_D+e_L+7=9\).
The retained set for the conductor-one target is \(\{1\}\) with
weight \(1\), so the function is \(G_1\) at every scale.
\begin{figure}[!htbp]
\centering
\begin{tikzpicture}[x=1mm,y=1mm,
  window step/.style={argument block,text width=42mm,minimum height=15mm}]
  \node[window step,fill=FieldTint] (requirements) at (0,0)
    {\hyperlink{proof.join-size-requirements}{\shortstack{CM fields and combined\\size requirements}}};
  \node[window step,fill=TraceTint] (lattice) at (55,0)
    {\hyperlink{proof.uniform-lattice-excess}{\shortstack{Lattice inequality\\for each translate}}};
  \node[window step,fill=AnalysisTint] (pointwise) at (110,0)
    {\hyperlink{proof.fourier-exponent}{\shortstack{Pointwise implication\\by Fourier inversion}}};
  \node[window step,fill=ChoiceTint] (choose) at (0,-24)
    {\hyperlink{proof.choose-compatible-family}{\shortstack{Fix the family and \(x_*\),\\then choose the interval}}};
  \node[window step,fill=ChoiceTint] (inputs) at (55,-24)
    {\hyperlink{proof.same-family-hypotheses}{\shortstack{Use \(D(u)\le E_x\) for all \(u\)\\when the lattice hypotheses\\hold at the chosen \(x\)}}};
  \node[window step,fill=AnalysisTint] (apply) at (110,-24)
    {\hyperlink{proof.apply-full-window}{\shortstack{\(G_1(x)\le x^{1/2+9\varepsilon}\)\\
       for all real \(x\in[Y,Y^{3/2}]\)}}};
  \draw[argument arrow] (requirements.south) -- (choose.north);
  \draw[argument arrow] (lattice.south) -- (inputs.north);
  \draw[argument arrow] (choose.east) -- (inputs.west);
  \draw[argument arrow] (inputs.east) -- (apply.west);
  \draw[argument arrow] (pointwise.south) -- (apply.north);
\end{tikzpicture}
\caption{Fix the target, the loss bounds \(e_D=2\) and \(e_L=0\),
  and a sufficiently small \(\varepsilon>0\). Choose one family
  \((\mathfrak F,C_H)\) meeting the joint size requirements, including
  \(C_H>0\). Fix its comparison constants and common threshold \(x_*\).
  For each \(Y_0\ge2\), choose \(Y\ge\max(Y_0,x_*)\) with
  \([Y,Y^{3/2}]\subseteq\operatorname{Win}(\mathfrak F)\).
  At each real scale in this interval, obtain the lattice bound after
  supplying its period, rank and height hypotheses at the tied indices
  \(i(u)\). The Artin and character data come from \(i_0\). The shared
  degrees, cutoff and centering term identify the sum used in the
  pointwise implication.}
\label{fig:family-window}
\end{figure}


\teamanchor{proof.apply-full-window}%
Given \(Y_0\ge2\), the window property gives \(Y\ge\max(Y_0,x_*)\)
with \([Y,Y^{3/2}]\subseteq\operatorname{Win}(\mathfrak F)\).
Every real \(x\) in the interval has \(x\ge x_*\), so the pointwise
implication applies at every scale in Figure~\ref{fig:family-window}.
\begin{equation}
\label{eq:prime-window}
 G_1(x)\le x^{1/2+9\varepsilon}
 \qquad\text{for every real }x\in[Y,Y^{3/2}].
\end{equation}
The same family and \(x_*\) work as \(Y_0\) varies. The precision
can be any fixed positive number below the minimum of the cutoffs.

\subsection{The centered zeta transform}

\teamanchor{proof.centered-mellin-residues}%
Write
\[
 \Phi_1(s)=\int_1^\infty G_1(x)x^{-s-1}\,dx .
\]
The finite ramp sum is continuous in \(x\). From \(\Lambda(n)\le\log n\),
its absolute value is at most \(x\sum_{n\le x}(\log n)/n\).
Including the centering terms gives \(G_1(x)=O(x(1+\log^2x))\), so
the integral converges absolutely on \(\Re s>1\).
The Euler product converges there and gives a convergent logarithm.
Differentiating locally uniformly yields
\[
 -\frac{\zeta'(s)}{\zeta(s)}=\sum_{n\ge1}\frac{\Lambda(n)}{n^s}.
\]
\cite[(25.2.11)]{dlmf}
For \(\Re s>1\) the transform of a ramp is
\[
 \int_n^\infty (x/n-1)x^{-s-1}\,dx
 =n^{-s}\left(\frac{1}{s-1}-\frac{1}{s}\right).
\]
The sum of the absolute integrals at \(\sigma=\Re s>1\) is bounded
by \((\sigma(\sigma-1))^{-1}\sum_n(\log n)n^{-\sigma}\).
Interchanging sum and integral and transforming the two centering
terms of \eqref{eq:principal-ramp} gives
\[
 \Phi_1(s)=\frac{-\zeta'(s)/\zeta(s)}{s(s-1)}
                  -\frac{1}{(s-1)^2}+\frac{1+\gamma}{s-1}.
\]
The continuation of \(\zeta\) is meromorphic with unique pole at \(1\)
and \(\zeta(1+z)=z^{-1}+\gamma+O(z)\).
\cite[\S25.2(i) and (25.2.4)--(25.2.5)]{dlmf}
Hence \(-\zeta'(1+z)/\zeta(1+z)=z^{-1}-\gamma+O(z)\).
Dividing by \(z(1+z)\) gives \(z^{-2}-(1+\gamma)z^{-1}+O(1)\),
which the centering cancels.

The Euler product gives nonvanishing for \(\Re s>1\).
On the line at \(1+it\), \(t\ne0\), use the nonnegative sum
\[
 -\Re\left(3\frac{\zeta'(\sigma)}{\zeta(\sigma)}
       +4\frac{\zeta'(\sigma+it)}{\zeta(\sigma+it)}
       +\frac{\zeta'(\sigma+2it)}{\zeta(\sigma+2it)}\right)
  =2\sum_{n\ge1}\frac{\Lambda(n)}{n^\sigma}
                 \bigl(1+\cos(t\log n)\bigr)^2
 \quad(\sigma>1).
\]
If \(1+it\) were a zero of multiplicity \(m\ge1\), the coefficient
of \((\sigma-1)^{-1}\) as \(\sigma\downarrow1\) would be
\(3-4m-m_2<0\), where \(m_2\ge0\) is the zero multiplicity at
\(1+2it\) (zero if the value is nonzero). This contradicts the
displayed inequality.
\cite[\S25.10(i)]{dlmf}
On the real segment \(0<\sigma<1\), pairing consecutive terms
in the alternating zeta series gives
\[
 (1-2^{1-\sigma})\zeta(\sigma)
 =\sum_{j\ge1}\bigl((2j-1)^{-\sigma}-(2j)^{-\sigma}\bigr)>0.
\]
The prefactor is negative, so \(\zeta(\sigma)<0\).
\cite[(25.2.3)]{dlmf}

It follows that \(\Phi_1\) is meromorphic on \(\Re s>0\) with
one simple pole at each zero \(\rho\) in the open strip and with
residue
\begin{equation}
\label{eq:zeta-window-residue}
 \operatorname{Res}_{s=\rho}\Phi_1(s)
   =-\frac{m_\rho}{\rho(\rho-1)}
   =\frac{m_\rho}{\rho(1-\rho)} .
\end{equation}
Indeed, the factor \((s-\rho)^{m_\rho}\) in \(\zeta(s)\) gives
residue \(-m_\rho\) in its negative logarithmic derivative,
and \(s(s-1)\) is nonzero at \(\rho\). These poles are nonreal,
come in conjugate pairs and are locally finite in \(\Re s>0\).

\subsection{One-sided windows and the critical line}

\teamanchor{proof.one-sided-window-criterion}%
Use the following simple-pole case of Kaczorowski and Pintz. Let \(h\)
be real and locally integrable on \((0,\infty)\), with Mellin integral
\(\int_0^\infty h(x)x^{-s-1}\,dx\) absolutely convergent on some closed
right half-plane. Suppose this transform has a meromorphic continuation
to a neighborhood of \(\Re s\ge\Theta-c_0\), for \(\Theta>0,\ c_0>0\),
holomorphic on \(\Re s>\Theta\). Suppose its poles in the closed
half-plane are simple, nonreal, locally finite and paired by conjugation,
with supremum of real parts \(\Theta\). For every \(0<\delta<1\), all
sufficiently distant intervals \([X^{1-\delta},X]\) contain both a
strictly positive and a strictly negative value of \(h\).

To apply Part~II, cut the half-plane leftward from each of these
poles. The transform restricts holomorphically to the cut domain,
and near each pole its expansion in (1.5) has exponent \(-1\),
logarithmic degree \(0\), and nonzero coefficient equal to the residue.
The pole supremum is the boundary of the half-plane of holomorphy
whether or not a pole attains it. Conclusion (1.7) gives the two
signs. The same case in Part~I uses zero logarithmic term.
\cite[(1.5) and (1.7)]{z6kaczorowski-pintz-ii}~\cite[Thm.~1 and (2.2)]{z6kaczorowski-pintz}

\teamanchor{proof.zero-precision}%
Suppose \(\zeta(\rho)=0\) with \(1/2<\beta=\Re\rho<1\).
Choose \(1/2<\theta<\beta\), then choose \(0<\varepsilon\) below all
the preceding cutoffs and with \(9\varepsilon\le\theta-1/2\).
The fixed-precision bound \eqref{eq:prime-window} gives arbitrarily
distant intervals \([Y,Y^{3/2}]\) on which \(G_1(x)\le x^\theta\).

Set \(h(x)=G_1(x)-x^\theta\) for \(x\ge1\) and \(h(x)=0\) for
\(0<x<1\). It is real and locally integrable, and its Mellin
integral converges absolutely on \(\Re s\ge2\). The transform there
is \(\Phi_1(s)-1/(s-\theta)\). Let \(\Theta\) be the supremum
of the real parts of the poles of \(\Phi_1\). By
\eqref{eq:zeta-window-residue}, \(\Theta\ge\beta>\theta\)
and \(\Theta\le1\). Choose \(c_0=(\Theta-\theta)/2>0\).
On the closed half-plane \(\Re s\ge\Theta-c_0\) the correction
term is holomorphic, and the poles are exactly those of \(\Phi_1\)
there. They satisfy all the sign-change hypotheses and still have
supremum \(\Theta\). Apply the theorem with \(\delta=1/3\) at
\(X=Y^{3/2}\), taking \(Y\) large enough to exceed its threshold
for this fixed \(h\) and \(\delta\). An interval with
\(X^{1-\delta}=Y\) must contain \(h>0\), contradicting \(h\le0\)
on the entire interval.

\teamanchor{proof.zeta-reflection}%
We have excluded all zeros to the right of the critical line.
For \(\Re s<1/2\), \(s\ne0\), use the functional equation
\[
 \zeta(s)=2^s\pi^{s-1}\sin(\pi s/2)\Gamma(1-s)\zeta(1-s).
\]
\cite[(25.4.2)]{dlmf}
The argument \(1-s\) lies on the zero-free side and differs from \(1\).
The gamma factor is holomorphic and nonzero there.
\cite[\S5.2(i)]{dlmf}
The sine factor vanishes precisely when \(s\) is a negative even
integer in this region. Also \(\zeta(0)=-1/2\).
\cite[(25.6.1)]{dlmf}
Thus \eqref{eq:prime-window} implies that every zero other than the
negative even integers lies on \(\Re s=1/2\), the classical
Riemann hypothesis.

\subsection{Extension to Dirichlet \texorpdfstring{\(L\)}{L}-functions}

For a fixed primitive target \(\psi\), assume that the arithmetic and
period comparisons hold on its selected family. The window argument
then gives
\begin{equation}
\label{eq:primitive-target-window}
 G_{\mathcal K}(x)\le x^{1/2+9\varepsilon}
 \qquad(x\in[Y,Y^{3/2}]),
\end{equation}
for arbitrarily large \(Y\) at each sufficiently small fixed
\(\varepsilon>0\). The retained \(\mathcal K\) consists of the
primitive associates of the powers of \(\psi\), with fixed positive
weights, and contains \(\psi\).
\releasecite{Development/AutomorphicArithmetic/PrimitiveCosets/LevelGmLeOnFamilyScales.lean}{28}{The fixed-target window comparison}
is applied to the primitive character class in the
\releasecite{Development/AutomorphicArithmetic/CharacterBoxes/ZeroFreeAllOfTopG.lean}{24}{Dirichlet construction}.

\paragraph{The Mellin poles.}
For any finite set \(\mathcal F\) of primitive Dirichlet characters
with weights \(a_\chi>0\), put
\[
 G_{\mathcal F}(x)=\sum_{\chi\in\mathcal F}a_\chi\Re G_\chi(x),
 \qquad
 \Phi(s)=\int_1^\infty G_{\mathcal F}(x)x^{-s-1}\,dx .
\]
The same absolute-convergence calculation as for \(\Phi_1\) works,
since \(|\chi(n)|\le1\). For primitive \(\chi\ne1\) its Euler
product, differentiated on \(\Re s>1\), gives
\(-L'(s,\chi)/L(s,\chi)=\sum_{n\ge1}\Lambda(n)\chi(n)n^{-s}\).
\cite[(25.15.1)--(25.15.2)]{dlmf}
Using the ramp transform and \eqref{eq:centred-prime-sums} yields
\[
 \Phi_\chi(s)=\frac{-L'(s,\chi)/L(s,\chi)}{s(s-1)}
                 +\frac{L'(1,\chi)/L(1,\chi)}{s-1}.
\]
Here \(L(s,\chi)\) is entire and \(L(1,\chi)\ne0\).
\cite[\S25.15(i), (25.15.9)]{dlmf}
The centering cancels the residue of the quotient at \(s=1\).
The Euler product gives nonvanishing for \(\Re s>1\), and
Montgomery and Vaughan's zero-free region gives
\(L(1+it,\chi)\ne0\) for all real \(t\) for nonprincipal
\(\chi\). The exceptional zero in that theorem is real and
strictly less than \(1\).
\cite[Thm.~11.3]{z6montgomery-vaughan}
Thus the poles of the transforms in \(\Re s>0\) lie in the open strip.

Conjugating the Dirichlet series and continuing gives
\(L(s,\bar\chi)=\overline{L(\bar s,\chi)}\), so
\(G_{\bar\chi}(x)=\overline{G_\chi(x)}\). The transform of
\(\Re G_\chi\) is \((\Phi_\chi+\Phi_{\bar\chi})/2\).
Write \(m_\chi(\rho)\) for the zero multiplicity at \(\rho\),
with value zero at nonzeros. The calculation of
\eqref{eq:zeta-window-residue} gives
\[
 \operatorname{Res}_{s=\rho}\Phi(s)
 =\frac{\sum_{\chi\in\mathcal F}
        a_\chi\bigl(m_\chi(\rho)+m_{\bar\chi}(\rho)\bigr)}
        {2\rho(1-\rho)}
 \qquad(0<\Re\rho<1).
\]
Thus a zero of any retained character gives a simple pole even
when shared by several characters. At a real pole the residue is
positive, since \(0<\rho<1\).

\paragraph{Positive majorants and real poles.}
The release uses a one-sided criterion for a real measurable \(g\)
on \([0,\infty)\) with \(|g(u)|\le M e^{2u}\). Assume its Laplace
transform \(\Phi_g\) extends meromorphically from \(\Re s>2\) to
\(\Re s>0\), with simple poles and positive residues at real poles.
Fix \(0<\theta<1\) and \(\lambda>1\). Assume that for every
\(\eta>0\) there are arbitrarily large \(U\) with
\begin{equation}
\label{eq:positive-majorant-window}
 \int_U^{\lambda U}g_+(u)e^{-\theta u}\,du\le e^{\eta U},
 \qquad g_+(u)=\max(g(u),0).
\end{equation}
Then \(\Phi_g\) has no pole in \(\Re s>\theta\).
The positive part is the least nonnegative majorant of \(g\).
ai agent proves the criterion by estimating derivatives at a real
center and a vertical translate, then comparing their growth at the
nearest pole.
\releasecite{Development/AutomorphicArithmetic/DedekindLFunctions/LandauTheorems/landauWindow.lean}{40}{The derivative proof in the release}

Here is the deduction from the published sign-change theorem.
If \(\Phi_g\) has a pole \(\rho\) with \(\Re\rho>\theta\), fix
\(\eta>0\) with \(\alpha=\theta+\eta<\Re\rho\).
Extend \(g\) by zero to \(u<0\). The average
\(\widetilde g(v)=\int_{v-1}^v g(u)\,du\) has Laplace transform
\[
 \widetilde\Phi(s)=\Phi_g(s)\frac{1-e^{-s}}{s}
\]
by Fubini on \(\Re s>2\). The multiplier is nonzero on \(\Re s>0\)
and positive for real \(s>0\). Thus the average preserves both the
poles there and the signs of real residues. On each supplied window,
for \(U+1\le v\le\lambda U\),
\[
 \widetilde g(v)\le
 e^{\theta v}\int_{v-1}^v g_+(u)e^{-\theta u}\,du
 \le e^{\theta v+\eta U}\le e^{\alpha v}.
\]
Choose \(1<\lambda'<\lambda\). For all sufficiently large \(U\),
\([V,\lambda' V]\subseteq[U+1,\lambda U]\) at \(V=U+1\).

The real poles \(B\) of \(\widetilde\Phi\) above \(\alpha\) are
finite by meromorphy around the compact segment \([\alpha,2]\)
and absolute convergence on \(\Re s>2\). Write \(b_\beta>0\) for
their residues. Choose \(t\in[1,2]\) such that none of the
\(\beta\pm it\), \(\beta\in B\), is a pole of \(\widetilde\Phi\).
Only finitely many values of \(t\) are excluded. Put
\[
 h(x)=\widetilde g(\log x)-x^\alpha
       -\sum_{\beta\in B}b_\beta x^\beta(1+\cos(t\log x))
 \qquad(x\ge1)
\]
and \(h(x)=0\) below \(1\). This function is real and locally
integrable and its Mellin integral converges absolutely on \(\Re s\ge3\).
The transform is
\[
 \widetilde\Phi(s)-\frac{1}{s-\alpha}
 -\sum_{\beta\in B}b_\beta
 \left(\frac1{s-\beta}+
       \frac{1}{2(s-\beta-it)}+\frac{1}{2(s-\beta+it)}\right).
\]
Each pole at \(\beta\in B\) is canceled and two simple conjugate
poles at \(\beta\pm it\) take its place, with residues \(-b_\beta/2\).
Thus the supremum \(\Theta>\alpha\) of the pole real parts stays
the same. For \(0<c_0<\Theta-\alpha\), the poles to the right of
\(\Theta-c_0\), including that line, are all nonreal and simple.
The transform satisfies the other sign-change hypotheses by
meromorphy on \(\Re s>0\) and conjugation. The correction is
nonnegative before subtraction, so \(h(x)\le0\) throughout
\([e^V,e^{\lambda' V}]\) for arbitrarily large \(V\).
With \(X=e^{\lambda' V}\), \(\delta=1-1/\lambda'\), this contradicts
the sign-change theorem.

\paragraph{Primitive and imprimitive targets.}
For a hypothetical zero \(\rho\) of \(L(s,\psi)\) with
\(1/2<\Re\rho<1\), choose \(1/2<\theta<\Re\rho\) and a fixed
sufficiently small \(\varepsilon>0\) with \(9\varepsilon\le\theta-1/2\).
Use \(\mathcal F=\mathcal K\) and put \(g(u)=G_{\mathcal K}(e^u)\).
The growth bound \(O(e^u(1+u^2))\) gives the required \(Me^{2u}\).
By \eqref{eq:primitive-target-window}, the integral in
\eqref{eq:positive-majorant-window} is at most \((\lambda-1)U\)
on arbitrarily distant windows at \(\lambda=3/2\).
For each \(\eta>0\) this is at most \(e^{\eta U}\) once \(U\)
is large. The Mellin residue computation and \(\psi\in\mathcal K\)
give a pole at \(\rho\), contradicting the criterion.

Let \(\chi\) be a primitive nonprincipal character of conductor \(d\)
and parity \(a\in\{0,1\}\), defined by \(\chi(-1)=(-1)^a\). Its
completed function is entire and satisfies
\begin{equation}
\label{eq:primitive-functional-equation}
 \begin{aligned}
 \Lambda_\chi(s)&=(d/\pi)^{(s+a)/2}\Gamma\bigl((s+a)/2\bigr)L(s,\chi),\\
 \Lambda_\chi(s)&=\epsilon_\chi\Lambda_{\bar\chi}(1-s)
       \qquad(s\in\mathbf C),
 \end{aligned}
\end{equation}
where \(\epsilon_\chi=i^{-a}d^{-1/2}\sum_{r\bmod d}\chi(r)e^{2\pi i r/d}\)
has modulus \(1\).
\cite[(10.17) and Cor.~10.8]{z6montgomery-vaughan}
The gamma and conductor factors are finite and nonzero in the open
strip, so a zero there at \(s\) gives one at \(1-s\) for \(\bar\chi\).
Thus the window bound \eqref{eq:primitive-target-window} for the two
fixed targets \(\chi\) and \(\bar\chi\) puts every zero in the open
strip on \(\Re s=1/2\).

For any Dirichlet character \(\chi\) modulo \(N\ge1\), let
\(\chi^*\) be its inducing primitive character of conductor \(d\mid N\).
The change-of-conductor formula is a meromorphic identity
\[
 L(s,\chi)=L(s,\chi^*)\prod_{\substack{p\mid N\\p\nmid d}}
                        (1-\chi^*(p)p^{-s}).
\]
\cite[(25.15.4)]{dlmf}
Each displayed Euler factor is holomorphic and nonzero in
\(\Re s>0\), since \(|\chi^*(p)p^{-s}|<1\) there.
For a principal character \(\chi^*=1\), so \(L(s,\chi)\) has a simple
pole at \(s=1\) of residue \(\prod_{p\mid N}(1-p^{-1})=\varphi(N)/N\).
Consequently, given the window bound for each primitive target,
every \(L(s,\chi)\) is nonzero in \(\Re s>1/2\) away from the
principal pole. Every zero in \(0<\Re s<1\) then has real part \(1/2\).
ai agent uses
\releasecite{Development/AutomorphicArithmetic/CharacterBoxes/Imprimitive.lean}{5}{the same change of conductor}
for the right-hand nonvanishing in the
\releasecite{Main.lean}{8}{formal Dirichlet theorem}.
\par
